\documentclass[lualatex,book,paper=a4paper,jafontsize=11.5pt,jafontscale=1.05]{jlreq}

%フォント:
\usepackage{luatexja-fontspec}
\usepackage{hack-fonts}

%表紙用:
\usepackage{hack-cover}

% リンク
\usepackage[hidelinks]{hyperref}

% 画像
\usepackage{graphicx}


% 単位 
\usepackage{siunitx}

\usepackage{float}    % [H] 指定を有効にするため必須
\usepackage{array}    % p{6cm} などの列幅指定の拡張
\usepackage{booktabs} % より美しい表の罫線（推奨、今回は \hline のままでも動きます）
\usepackage{tabularx} % ★tabularx環境とX列を使うために必須！
\usepackage[table]{xcolor}

\title{ハッカソン}
\subtitle{～計画書と設計書～}
\team{チーム40}
\author{
        24G1021 : 宇井　祐真\\
        24G1026 : 梅野　大誠\\
        24G1136 : 山口　侑希\\
        25G1020 : 薄井　悠希\\
        25G1041 : 菊池　侑人\\
        25G1059 : 佐藤　涼菜 
}
\date{\today}
\version{1.0}

%%%%%%%%%%%%%%%%%%%%%%%%%%%%%%%%%%%%%%%%%%%%%%%%
\begin{document}
\maketitle

\setcounter{tocdepth}{2}
\tableofcontents \thispagestyle{empty}

%%%%%%%%%%%%%%%%%%%%%%%%
%%%%%%%%%%%%%%%%%%%%%%%%
%%%%%%%%%%%%%%%%%%%%%%%%
\chapter{プロジェクトの計画書} \label{chap:project-plan}\setcounter{page}{1}
%1%%%%%%%%%%%%%%%%%%%%%%
\section{プロジェクトの概要}
%%%%%%%%%%%%%%%%%%%%%%%%
\subsection{プロジェクトの題目}
マジカル☆観覧車
%%%%%%%%%%%%%%%%%%%%%%%%
\subsection{プロジェクトの目的}
% プロジェクトの目的と最終的なゴール，解決すべき課題を記述する．
近年，私たちは動画サイトやストリーミングサービスによって，誰でも手軽に音楽コンテンツを楽しむことができるようになった．
特に，日本独自の文化として，ゲームやアニメといったサブカルチャー発の音楽コンテンツが盛んである．
このようなコンテンツは，上記のようなデジタル技術の発展の恩恵を受ける一方で，
実際の体験に重きを置くライブイベントやコンサートなどとの相乗効果を得て，市場規模の拡大に大きく貢献している\cite{cul}\cite{live}．
実際に，一般社団法人コンサートプロモーターズ協会が行ったライブ市場調査によると，2025年にはライブの年間売上として6443億円を記録しており，
入場者数は5999万人を記録した\cite{acpc}．
これらのイベントが市場シェアを拡大する理由として，
ライブ演出における音響や照明，身体で感じる振動などは，その演奏状態を複数人へ同時に共有可能であり，
空間全体としての一体感を形成する役割を持っている点が挙げられる．
つまり，ユーザが音楽を聴くだけにとどまらず，音響や照明などの物理的，視覚的な技術をその場で一体となって
楽しむといった，没入感の高い体験を得られるという点が理由の一つである．\par
ここで，没入感はどのように創出されるのだろうか．
Natalia Cooperら(2018)は，VR環境においてのシミュレーションの充実度を高めるため，
再現の難しい感覚を別感覚による代替フィードバックとして与える実験を行った．
その結果，VR上の視覚的な刺激のみよりも，音や触覚といった異なる刺激を追加することが，
ユーザのシステムに対する没入感・関与感を向上させることを明らかにした\cite{Natalia2018}．
このことから，複数の刺激による多感覚的な情報の提示は，ユーザ体験における没入感の向上に有効であるとわかる．\par
また，HCI分野において，ユーザの主体的な動作がシステムに対する没入感を高めることができることが示されている．
Martinら(2024)によると，ライブコンサート環境において，観客は受動的な聴取をするのでなく，
能動的な参加を行うことが，ライブの活気や観客の感情性を形成すると結論づけた\cite{Martin}．
よって，能動的な活動参加もまた，システムの没入感を高める要因となるとわかる．\par
一方，現状の自動演奏技術分野については，デジタル技術の発達により高精度な同期や楽曲再現を行うことができる．
例えば，yamahaよりリリースされているDisklavier ENSPIREシリーズは，
光学センサやサーボモータを用いることで，演奏情報の高精度再現を行うことが可能であり，
ユーザに対して，物理的な鍵盤の動きを用いての視覚的フィードバックも可能としている\cite{yamaha}．
しかしながら，このフィードバックは多くの場合，利用者１人に対してのものとなり，
空間全体を考えると，他に与える刺激は限定的になる場合がある．
よって，現在の自動演奏技術分野は，視覚的なフィードバックを追加し没入感を高めているものの，
その実，空間的フィードバックは限定的であると言える．\par
従って，本プロジェクトではこの問題を解決するため，
自動演奏技術に空間的フィードバックを追加した楽器間通信システムであるマジカル☆観覧車を提案する．

%%%%%%%%%%%%%%%%%%%%%%%%
\subsection{既存の技術とプロジェクトの独自性}
本節では，プロジェクトの設計計画にあたり，既存技術の分析を通して本プロジェクトの独自性について検討する．
まず，音楽に視覚的なフィードバックを追加した例として，DMX512信号(以下，DMX信号と示す)を用いたシステムが挙げられる．
DMX信号とは，操作卓と調光機の制御に用いられ，現在では舞台やコンサートなどの場面で演出に合わせた光量制御に用いられている\cite{sekine2021}．
このシステムは，音楽だけでなく光といった異なる刺激を与えることで，没入感を高めることができる．
しかしながら，利用されるシチュエーションを考えると，システムの対象である観客は現象を受動的に認識する場合が多く，観客が演奏に参加できない場合が考えられる．
ここで，本提案では空間フィードバックの一環として指揮デバイスを導入している．
このデバイスを用いることで，ユーザが主体的に演奏に参加することができ，身体的な動作関与による没入感の向上を期待することができると考えられる．\par
また，物理的な操作による音楽の制御を行うシステムの例として，ReacTableが挙げられる．
ReacTableとは，卓型のディスプレイと複数の物理モジュールによって構成されるシンセサイザの一種で，
身体的な物理動作を通して演奏の制御，ディスプレイによる視覚的フィードバックが可能になる\cite{reactable}．
しかしながら，ディスプレイ内に視覚的なフィードバックが行われる関係上，
空間全体へのフィードバックを与えるインターフェースとして扱われてはいない．
ここで，本提案では，指揮デバイスに加え観覧車型デバイスを導入している．このデバイスを用いることで，ユーザ以外からも演奏状態の把握や共有ができ，
空間全体への多感覚的フィードバックを可能にすることが期待できると考えられる．\par
以上より，本提案は既存研究の空間共有性と主体的な演奏制御技術を組み合わせることで，
空間全体への感覚フィードバックを行い，広く没入感や関与感を高めることができると考えられる．


%%%%%%%%%%%%%%%%%%%%%%%%
%%%%%%%%%%%%%%%%%%%%%%%%
\section{スコープ・マネジメント}
%%%%%%%%%%%%%%%%%%%%%%%%
\subsection{成果物}
本プロジェクトで製作する成果物は，同期開始に指揮棒，中継機に観覧車を用いた楽器演奏システムである．
指揮動作によって中継機と楽器デバイスの開始と終了，BPMと音量情報を送信する指揮デバイス，
指揮デバイスから情報を受け取り，BPMに応じて回転する観覧車型中継機デバイス，
光信号によってBPMを感知し，その速さで音楽を再生する楽器デバイスの３つに分割し製作する．楽器デバイスは，Arduinoで楽譜情報を作成し，Processingから楽器音声を再生する動作と，BPM変化や音声変更の処理を行い演奏情報を変更するシステムを制作する．
また，中継機ではLEDMatrixに現在のBPM数を表示するように製作を行う．
そして，3Dプリンタを用いた指揮デバイス機構，観覧車機構，レーザー光送信機構，
Arduinoを用いたモータ制御機構を成果物とする．

%%%%%%%%%%%%%%%%%%%%%%%%
\subsection{プロジェクトの対象範囲と非対象範囲}
本プロジェクトの対象範囲は，センサーやUDPを用いた無線通信，processingを用いた楽器演奏，3Dプリンタを用いた観覧車製作や中継機デバイス，指揮デバイスのハードウェア製作である．非対象範囲は，有線による情報転送，Arduinoからの音源再生，信号の暗号化である．

%%%%%%%%%%%%%%%%%%%%%%%%
\subsection{機能分解構造FBS}
機能分解構造FBSは，以下の図\ref{fbs}の通りである．観覧車型中継機デバイスを用いた楽器間通信デバイスの作成にあたって，ユーザ視点の機能として三つに分けられる．一つ目は楽器の演奏，二つ目に通信同期関係，三つ目に見た目の華やかさである．その三要素を詳細に分割すると，楽器演奏に関しては多彩の楽器の音色や高度な音色再現が挙げられる．通信同期に関しては，ズレの少ない同期や変化する速度が挙げられ，見た目の華やかさに関しては，視覚的に伝わりやすい表示や見て興味を誘うような動作が挙げられる．
\begin{figure}[h]
\centering
 \centering
\includegraphics[width=15cm]{./figs/FBS.png}
\caption{FBS図}
\label{fbs}
\end{figure}

%%%%%%%%%%%%%%%%%%%%%%%%
\subsection{製品分解構造PBS}
次に製品分解構造PBSについて説明する．PBSは以下の図\ref{pbs}の通りである．観覧車型中継機デバイスを用いた楽器間通信デバイスをサブ成果物に分解すると，Processing，楽器，楽曲表現，通信の四つに分割される．Processingに関しては，音声波形の生成と音声の再生に分割され，音声波形の生成は基礎波形情報と適切なエフェクトが成果物としてあげられ，音声の再生は発音機構が成果物として挙げられる．楽器に関しては，同期通信の開始信号が必要であり，そこから照度センサと開始信号からのBPM計算処理プログラムが成果物で挙げられる．楽曲表現に関しては，演奏情報や回転デバイスに挙げられ，それぞれLEDマトリクスとモーター，観覧車本体が成果物として挙げられる．通信に関しては，指揮デバイスに関する通信，中継機に関する通信，同期通信の三つに分割される．指揮デバイスについては加速度の設計と演奏情報の更新機構が成果物としてあげられ，中継機については演奏情報の一斉送信のプログラム設計，同期通信については光学設計とモータ回転速度の設計が成果物としてあげられる．

\begin{figure}[h]
\centering
 \centering
\includegraphics[width=17cm]{./figs/PBS.png}
\caption{PBS図}
\label{pbs}
\end{figure}
%%%%%%%%%%%%%%%%%%%%%%%%

\subsection{FBSとPBSの対応関係}
FBSとPBSの対応関係について説明する．関係図は以下の図\ref{fp}の通りである．矢印の左側がFBS要素であり，右側がPBS要素である．多様な楽器の音色の要素には基礎波形情報と，音色を鳴らすための発音機構が対応する．高度な音色再現には適切なエフェクトが対応し，視覚的に伝わりやすい表示にはLEDマトリクスが対応する．
ズレの少ない同期では，演奏情報の一斉送信や光学設計，BPM計算処理や照度センサの成果物が対応する．速度変化が可能には，モータ回路や回転速度の設計，演奏情報の更新機構が対応し，興味を誘う動作には加速度センサや観覧車本体の成果物が対応する．
\begin{figure}[h]
\centering
 \centering
\includegraphics[width=15cm]{./figs/fbspbs.png}
\caption{FBSとPBSの対応関係図}
\label{fp}
\end{figure}

%%%%%%%%%%%%%%%%%%%%%%%%
%%%%%%%%%%%%%%%%%%%%%%%%
\section{作業計画}
%%%%%%%%%%%%%%%%%%%%%%%%
\subsection{作業分解構造WBSと管理番号}
作業分解構造WBSと管理番号について説明する．WBSは設計フェーズ，実装フェーズ，テストフェーズに分けて説明する．それぞれ図\ref{k}，\ref{j}，\ref{t}の通りである．図\ref{k}について，1100は基本設計についてであり，指揮デバイス，中継機デバイス，楽器デバイスの設計をそれぞれ割り当てる．1200では詳細設計を割りあて，中継機デバイス，指揮デバイス，楽器，通信タイミング制御のハードウェア・ソフトウェア面の内容を作成する．図\ref{j}について，2100はソフトウェア，2200はハードウェアとし，実際に作業する内容を詳細に分割する．図\ref{t}について，実装フェーズで作成したものを単体でテストする内容を3100に，そしてそれらを結合して行うテストを3200に分割する．

\begin{figure}[h]
\centering
 \centering
\includegraphics[width=17cm]{./figs/wbsk.png}
\caption{設計フェーズのWBS}
\label{k}
\end{figure}

\begin{figure}[h]
\centering
 \centering
\includegraphics[width=17cm]{./figs/wbsj.png}
\caption{実装フェーズのWBS}
\label{j}
\end{figure}

\begin{figure}[h]
\centering
 \centering
\includegraphics[width=17cm]{./figs/wbst.png}
\caption{テストフェーズのWBS}
\label{t}
\end{figure}

%%%%%%%%%%%%%%%%%%%%%%%%
\subsection{アローダイアグラムとクリティカルパス}
成果物の完成までのアローダイアグラムは図\ref{aro}の通りである．
実装フェーズからテストフェーズまで消費する日付を考慮してつなげた結果，クリティカルパスは2131，2132，2135，3141，2224，3121，3211，3212，3231，3242という順序となり，日数は33日となった．この日数は制作期間である40日を下回る日数であるため，妥当な計画だとわかる．
\begin{figure}[h]
\centering
 \centering
\includegraphics[width=17cm]{./figs/アロー.png}
\caption{アローダイアグラム}
\label{aro}
\end{figure}

%%%%%%%%%%%%%%%%%%%%%%%%
\subsection{役割分担}
WBSを基に役割分担を作成する．役割分担表は，以下の\ref{yaku}の通りである．
分担の振り分けは主に楽曲表現型，楽曲系，通信系に分担し，それぞれハッカソン1と2でペアとなり製作を行う．
楽曲表現系は山口と佐藤，楽曲系は梅野と薄井，通信系は宇井と菊池が行う．

\begin{figure}[h]
\centering
 \centering
\includegraphics[width=17cm]{./figs/WBS.png}
\caption{WBSを基にした役割分担表}
\label{yaku}
\end{figure}

%%%%%%%%%%%%%%%%%%%%%%%%
\subsection{ガントチャート}
ガントチャートは図\ref{jgan}，\ref{tgan}の通りである．図\ref{jgan}は実装フェーズのガントチャートであり，三週目の途中までにそれぞれの単体作業を完成させる流れとなっている．図\ref{tgan}はテストフェーズのガントチャートであり，単体のテストはそれぞれの項目が完成次第テストを行って確認し，その後三週目以降から結合テストで動作するのを確認するという計画となっている．四週目終わりまでにすべての結合テストを実施し，最後の五週目は予備日として間に合わなかった際に予備日を用いて結合テストまでを行えるようにする．

\begin{figure}[h]
\centering
 \centering
\includegraphics[width=18cm]{./figs/gant1.png}
\caption{実装フェーズのガントチャート}
\label{jgan}
\end{figure}

\begin{figure}[h]
\centering
 \centering
\includegraphics[width=18cm]{./figs/gant2.png}
\caption{テストフェーズのガントチャート}
\label{tgan}
\end{figure}

%%%%%%%%%%%%%%%%%%%%%%%%
%%%%%%%%%%%%%%%%%%%%%%%%
\clearpage
\section{検証と妥当性確認: V\&V計画}
%%%%%%%%%%%%%%%%%%%%%%%%
\subsection{プロジェクトの成果物に対する有効指標MOE}
本プロジェクトの成果物に対する効果指標MOE(Measure of Effectiveness)は以下の5点とする．
作成するシステムは，「指揮デバイスと観覧車を用いた，空間的フィードバックにより高い没入感を実現する」という目的をもつため，人が知覚や体験が期待通りの効果を発揮
するかを評価軸とする．
\begin{itemize}
	\item 輪唱のタイミングに対して不快に感じる程度の遅延がないか
	\item ユーザの操作が正しく反映されるか
	\item 音色の再現ができているか
	\item BPMの回転速度がずれていないか
	\item エンターテイメント性があるか
\end{itemize}

%%%%%%%%%%%%%%%%%%%%%%%%
\subsection{有効指標MOEに対応する性能指標MOP}
効果指標MOEに対応する性能指標MOPとして，輪唱のタイミングに対して不快に感じる程度の遅延がないという基準を遅延時間が50ms未満とする．
50 ms未満という数値の決定するにあたり，ハース効果を根拠として採用する．
ハース効果とは，類似した二つの音が時間差で知覚された際，その時間差が一定時間内であれば，
聴覚がそれらを一つの音として知覚し，先に到達した音源から発生した音だと判断する現象である
\cite{bib:haas}．
そのため，作成するシステムの輪唱において，それぞれの楽器の演奏に50msを超えるズレが
発生すると人は演奏がずれていると知覚すると考えられる．
そのため，本システムの輪唱の遅延時間は50ms未満を技術的目標とする．

ユーザの操作が正しく反映されるかを評価するMOPとして，99\%の処理成功率を目指す．
ここで処理成功率の目標を99\%(ツーナイン)に設定した根拠は，一般的なwebサービスなどの開発環境に
おいて目標とされる数値のためである\cite{bib:twonine}．
本プロジェクトで作成するシステムは，ブレッドボード上で稼働する一つのシステムであるため，テスト環境での運用に近い．
そのため，ツーナインを参考に，作成するシステムの全体の処理成功率は99\%を目標とする．

音色の再現ができているかを評価するMOPとして，複数人に作成した音色を聞いてもらい主観評価で再現ができていると評価されることを目指す．
具体的な評価方法としては，複数人に対して，実際の楽器音と作成した楽器音を聴かせて，それらがどれだけ似ていたかを以下に示す1〜5で回答してもらい評価する．
\begin{itemize}
	\item 1:全く似ていない
	\item 2:あまり似ていない
	\item 3:どちらともいえない
	\item 4:ある程度似ている
	\item 5:非常に似ている
\end{itemize}
提示するサンプルは，順序による先入観を排除するため，
本プロジェクトで作成する楽器音と実際の楽器音をランダムな順序で提示する．
また，比較する7種類の楽器(ピアノ，ストリングス4種，フルート，ドラム)の音色は，それぞれの特性が現れる高音域，中音域，低音域の3種類で評価する．
これにより，広い音域での再現性を評価する．
ただし，ドラムの音色は1種類で評価する．
具体的な目標として，それぞれの音色の回答結果の中央値と最頻値が似ているとされる4.0以上と設定する．

BPMの回転速度がズレていないかに対応するMOPとして，回転速度とBPMのズレが\pm 1.5\%以内になることを目指す．
なぜならば，人が知覚できるテンポの変化の最小は1.5\%程度であり，それより小さくするためである\cite{bib:tempo}．
計測方法として，まずそれぞれのBPMで回転している観覧車をスマートフォンのカメラで動画撮影し，一周するまでに何枚のフレームがあるかを確認する．
次に，(\ref{eq:roll})式で一周にかかる時間を計測する．
最後に，それと(\ref{eq:bpmtoroll})式で表されるBPMから算出される理想的な一周にかかる時間との誤差率を(\ref{eq:riso})式を用いて算出する．
\begin{equation}
	\mbox{一周にかかる時間 (s)} = \frac{\mbox{フレーム数}}{\mbox{カメラのfps}}
        \label{eq:roll}
\end{equation}

\begin{equation}
	\mbox{一周にかかる理想時間 (s)} = \frac{60}{\mbox{BPM}} \times \mbox{4小節にかかる拍数}
	\label{eq:bpmtoroll}
\end{equation}

\begin{equation}
	\mbox{誤差率 (\%)} = \left| \frac{\mbox{一周にかかった時間} - \mbox{一周にかかる理想時間}}{\mbox{一周にかかる理想時間}} \right| \times 100
	\label{eq:riso}
\end{equation}

エンターテイメント性があるかに対応するMOPとして，複数人にアンケートを取り，顧客満足度CSATで75\%以上の範囲を目指す．
本プロジェクトで作成するシステムでは，エンターテイメント性を評価するために一般的な顧客満足度CSAT(Customer Satisfaction Score)を導入して，独自の評価軸に基づいたアンケート調査を実施する．
エンターテイメント性の評価をするため，以下に示す質問を設計した．各質問を5段階(1:大変不満，2:不満，3:どちらともいえない，4:満足，5:大変満足)のリッカート尺度で回答させる．
\begin{itemize}
	\item 観覧車の回転と楽器が演奏をしている様子を見て，空間的な楽しさや驚きに満足したか？
	\item 指揮棒を振ったり，可変抵抗器を操作して，演奏のテンポなどがリアルタイムに変化する操作感に満足したか？
	\item 光と音の動きが一体になったアトラクションとして，総合的な楽しさに満足したか？
\end{itemize}
アンケートの回答の評価としては，各質問において4(満足)または，5(大変満足)と回答した
回答者の割合を算出し，これを各項目のCSATスコア(\%)とする．CSATスコアの計算式は(\ref{eq:CSAT})式の通りである．
\begin{equation}
	CSATスコア = \left(\frac{4または5の回答者数}{回答者数}\right) \times 100 \label{eq:CSAT}
\end{equation}
算出されたCSATスコアは，一般的な企業のスコアの平均が76.7\%であることから，この数値と同等の75\%を目標値として設定する\cite{bib:ACSI}．




%%%%%%%%%%%%%%%%%%%%%%%%
\subsection{システムテストで評価する性能指標MOP}
システムテストフェーズでは，すべてのコンポーネントを統合した「システム全体」の完成状態で，最終的なMOPおよびMOEの評価実験を実施する．評価項目および具体的な手法は以下の通りである．

\begin{itemize}
	\item 輪唱タイミングの遅延計測(MOP目標：50ms未満)
	\\ 演奏開始信号が送信されてから各楽器が音を出すまでの時間，および楽器間の演奏ズレをmillis関数などを用いてミリ秒単位で計測し，ハース効果の許容限界である50ms未満を満たしているかを検証する．
	
\item 操作の反映度およびエンターテインメント性の主観評価(MOP目標：処理成功率99.0\%以上，CSATスコア75.0\%以上)
	\\ 開発メンバー以外の第三者を被験者（複数人）として，実際にシステムを操作してもらうユーザーテストを実施する．操作ログから入力に対するシステムの処理成功率を算出すると同時に，体験後に5段階リッカート尺度のアンケート調査を行い，式(\ref{eq:CSAT})を用いて算出されたCSATスコアが目標値に達しているかを評価する．
\end{itemize}
%%%%%%%%%%%%%%%%%%%%%%%%
\subsection{結合テストで評価する性能指標MOP}
結合テストフェーズでは，観覧車本体のハードウェア実装と制御プログラム，およびBPM指令を行う通信部を結合した段階で，主に物理的な同期性能の検証を行う．

\begin{itemize}
	\item 回転速度とBPMのズレに関する評価(MOP目標誤差率$\pm$1.5\%以内)
	\\ 指揮デバイスから様々なBPM（低速・中速・高速）の指令を送信し，それぞれの速度で回転する観覧車をスマートフォンのカメラで動画撮影する．
	\\ 撮影した動画から一周に要した「フレーム数」を目視または画像解析でカウントし，式(\ref{eq:roll})，式(\ref{eq:bpmtoroll})，および式(\ref{eq:riso})を用いて実際の誤差率を算出する．この誤差率が，人間のテンポ知覚限界である$\pm$1.5\%以内に収まっているかを検証する．
\end{itemize}
%%%%%%%%%%%%%%%%%%%%%%%%
\subsection{単体テストで評価する性能指標MOP}
単体テストフェーズでは，システムを結合する前の，各要素単体での基本性能を個別に評価する．

\begin{itemize}
	\item 音色の再現度評価(MOP目標：中央値および最頻値が4.0以上)
		\\ 自作した音源（ピアノ，ストリングス4種，フルート，ドラム）の単体出力音を録音する．複数人の被験者に対し，実際の楽器音と自作音をランダムな順序で聴取させる．
		\\ 各楽器の特性が現れる3つの音域(高・中・低音域，ドラムは1種類)について，1〜5の5段階評価で回答を得る．集計データの中央値および最頻値が「ある程度似ている」に相当する4.0以上を達成しているかを評価する．
	
	\item 各種センサー・入力デバイスの基礎精度検証
	\\ 指揮棒の判定プログラム，可変抵抗器のアナログ読み取り値，および照度センサによるLED読み取りを，シリアルモニタ等を用いて単体で計測し，処理精度を確認する．
\end{itemize}

%%%%%%%%%%%%%%%%%%%%%%%%
\subsection{プロジェクト中に監視する技術性能指標TPM}
プロジェクトの実行中に監視する技術性能指標TPM(Technical Performance Measure)として，
以下の5項目を定義する．これらは本システムの核心である「演奏情報通信機能」「観覧車の同期動作」「音響制御」の技術的リスクを追跡するための定量的な指標である．

\begin{itemize}
	\item \textbf{1. 入力・判定精度(\%)}
		\\ 指揮棒を振った際の判定精度，およびBPM・音量変更用可変抵抗器や照度センサの計測精度をそれぞれ計測する．
	
	\item \textbf{2. 通信・制御遅延時間(ms)とパケット到達精度(\%)}
		\\ 指揮棒操作や音量変更，演奏開始信号が送信されてから，観覧車および各楽器用のArduinoで受信・反映されるまでの遅延時間(ms)とデータ受信成功率(\%)を計測する．
	
	\item \textbf{3. モータ回転速度とBPMの同期精度(ms)}
	\\ 指揮デバイスで設定した目標BPMに対し，実際のモータ回転速度とのズレ時間を計測する．
	
	\item \textbf{4. 輪唱演奏の同期時間差(ms)}
	\\ 演奏開始信号が送信されてから楽器が実際に音を出すまでの時間，および各楽器が演奏を開始した時間と，輪唱として理想的な開始時間との時間差を計測する．
	
\item \textbf{5. 音色の再現度(類似度\%)}
	\\ 自作した音色が，自作の音に対して実際の楽器音との類似度の主観評価を5段階で計測する．
\end{itemize}

これらの定量的評価を毎週の授業での報告会で計測することで，
技術的リスクの早期発見，定量的な進捗管理を行う．
これにより，最終的なシステム性能要求MOPの達成を確実なものとする．

%%%%%%%%%%%%%%%%%%%%%%%%
\subsection{プロジェクト中に監視する指標TRL}
本プロジェクトの開発プロセスにおいて，技術的リスクの早期発見と進捗状況
の客観的な可視化を目的として，技術成熟度レベルTRL(Technology Readiness Level)を
導入する．
各開発フェーズにおける技術の成熟度は，以下の4段階に基づいて監視を行う．
\begin{itemize}
	\item 1.概念レベル：アイデアや原理が確立されており，説明できるレベル
	\item 2.基本動作レベル：単体で動作確認が可能なレベル
	\item 3.結合動作レベル：機能を結合し，実際の使用環境に近い環境で動作確認が可能なレベル
	\item 4.実運用レベル：実際の使用環境で安定的な稼働が確認できるレベル
\end{itemize}
本プロジェクトでは，毎週の授業中の報告会で各自が担当するシステムの構成要素が現在どの段階
にあるかを報告し，進捗状況を管理する．
これにより，主観的な進捗報告を排除し，定量的にプロジェクトを監視する．

%%%%%%%%%%%%%%%%%%%%%%%%
%%%%%%%%%%%%%%%%%%%%%%%%
\section{資源・調達管理}
資源・調達管理について説明する．今回作成するにあたってハードウェアで使用する電子部品は貸し出し機材と各自で購入し，作業開始日の5月26日までにすべて必要な機材を揃える．また，観覧車型中継機デバイスや指揮棒デバイスの工作は3Dプリンタを用いて製作を行う．この工作のための機材も作業開始日の5月26日までに作成し，開始日にすぐ作成に移れるような状態にする．

%%%%%%%%%%%%%%%%%%%%%%%%
%%%%%%%%%%%%%%%%%%%%%%%%
\section{リスク管理}
リスク管理について説明する．今回製作にあたって想定されるリスクとして，以下を挙げる．
\begin{itemize}
\item 想定より実装に時間がかかる
\item ハードウェアがうまく動作しない
\item  メンバー間で進捗に差が出る
\item グループ内の情報の伝達不備
\end{itemize}
それぞれのリスクに対して，以下の対策を考える．
\begin{itemize}
\item 実装の際にバッファを設け，時間がかかることを前提としたマネジメントを行
う．
\item データシート等から使用する素子の電気的特性を把握しておく．
\item 進捗が捗っているメンバーが，手助けを行う．また，確認した段階で現在の
進捗と実行プロジェクトを見直す．
\item 連絡を密に行い，認識のすり合わせを行う．また，作業者自身に説明させるこ
とで，認識を確認する．
\end{itemize}



%2%%%%%%%%%%%%%%%%%%%%%%
%%%%%%%%%%%%%%%%%%%%%%%%ここからシステム
%%%%%%%%%%%%%%%%%%%%%%%%
\chapter{システムの基本設計}
%%%%%%%%%%%%%%%%%%%%%%%%
\section{機能一覧}
%機能分解構造FBSの要素の説明
今回提案するものは，一つ目に信号を送る指揮デバイスの設計，二つ目に指揮デバイスから信号を受け取って動作する中継機デバイス，三つ目に楽器デバイスの設計である．
本節では指揮デバイスと中継機デバイス，楽器デバイスの機能について説明する．

始めに指揮デバイスについて説明する．
指揮デバイスとは演奏全体を統括するデバイスであり，主に以下の三つの機能を有する．
第一の機能は，演奏開始トリガーの送信である．
ユーザが指揮デバイスに取り付けられたボタンを押下すると，観覧車デバイスと楽器デバイスへ演奏開始トリガーを送信する．

%ユーザが指揮デバイスを上下に振ると，内蔵の加速度センサが振り動作に伴う加速度を検知する．そして，この値が設定した閾値を超過した瞬間に，観覧車デバイスと楽器デバイスへ演奏情報を送信する．

第二の機能は，楽曲のBPM変更である．ユーザが可変抵抗器を操作した後，指揮棒を振ると新しい設定BPM情報が中継機デバイスである観覧車へ送信される．これにより，観覧車のモータ回転速度がこの値に追従して変化する．そして，その回転に伴うレーザーの通過光を介して各楽器デバイスに情報が伝達され，リアルタイムに楽曲のBPMが変更および同期される．

第三に，楽曲の音量変更である．ユーザが可変抵抗器を操作した後，指揮棒を振ると音量情報が各楽器デバイスへマルチキャスト送信されることで楽曲の音量がリアルタイムに変更される．

次に，中継機デバイスについて説明する．
中継機デバイスとは，指揮デバイスからの情報を受け取り観覧車を作動させ，現在のBPMを表示するデバイスであり，主に以下の二つの機能を有する．
第一の機能は，上記で書いたBPMに合わせて回転を行い，楽曲のBPMを制御する機能である．観覧車に設置してあるレーザー受光のテンポでBPMを管理し，楽曲のテンポを視覚的にも理解できる機能である．
第二に，LEDMatrixの表示である．ここに指揮デバイスから受け取ったBPMを表示し，視覚的に現在流れている曲のテンポを理解できるような機能である．LEDMatrixはArduinoに付属したもので表示する．

最後に，楽器デバイスについて説明する．
楽器デバイスは観覧車デバイスの回転状態を基準として，自律的に演奏を行うデバイスであり，主に以下の四つの機能を有する．
第一の機能は，演奏開始の検知である．指揮デバイスが振り動作を検知すると，無線通信で楽器デバイスへと演奏開始情報が送信され楽曲の再生を開始する．
第二の機能は，楽曲のBPMリアルタイム同期である．演奏開始後もレーザー光の通過間隔を常時計測し，その間隔から現在のBPMを逐次算出する．
第三の機能は，楽器の音色再生である．ArduinoとUSBケーブルで接続されたPC上のProcessingを用いることで，各楽器の音色を出力し，演奏を再現する．
第四に，楽曲の音量リアルタイム反映である．指揮デバイスから受け取った音量情報をProcessingへ送信し，再生中の出力音量へ反映する．







%%%%%%%%%%%%%%%%%%%%%%%%
\section{システム構成図}
%ハードウェア，ソフトウェアなどの構成図を記載する．通信のつながりのようなもの
\subsection{システム全体の構成図}
今回提案する成果物全体の構成図は，図\ref{zentai}の通りである．指揮棒の振動を加速度センサーで感知した際に，Arduino1から無線で演奏情報を送信する．その後，情報を受け取ったArduino2が中継機である観覧車デバイスを作動させる．演奏の開始には指揮デバイスのボタンが押下された時，Arduino3〜6の楽曲演奏を開始する．観覧車デバイスが作動した時に，指揮デバイスから受け取った情報を基にBPMをLEDMatrixで表示を行う．観覧車が作動した時，
観覧車の8箇所に設置しているレーザーが楽器演奏デバイスにあるcdsセルに光を送信し，光を受信すると楽曲がPCのProcessingから流れる仕組みである．そして，流れてくるレーザーのテンポに従ってBPMが変化すると言う仕組みである．
\begin{figure}[h]
\centering
 \centering
\includegraphics[width=15cm]{./figs/zentai.png}
\caption{全体の構成図}
\label{zentai}
\end{figure}




\subsection{指揮デバイスの構成図}
指揮デバイスのシステム構成図を図\ref{fig:sys_siki}に示す．本デバイスは内蔵の加速度センサによりデバイスの振り動作を検知すると，観覧車デバイスと楽器デバイスにそれぞれに演奏情報が送信され，観覧車デバイスの回転が開始される．演奏の開始は，指揮デバイスのボタンを押すことで行う．また，ユーザによる操作により楽曲のBPMと音量を変更できる．
% また，ユーザによる操作により楽曲のBPMを変更することが可能であり，その情報は観覧車デバイスへ無線送信され，観覧車デバイスの回転を介して各楽器デバイスへと伝達される．
\begin{figure}[h] % h:ここ, t:上, b:下, p:独立ページ
        \centering % 中央揃え
        \includegraphics[width=0.8\columnwidth]{figs/siki_sys3.pdf} % 横幅を列幅の80%に指定
        \caption{指揮デバイスのシステム構成図} % 図の説明
        \label{fig:sys_siki} % 本文中で参照するためのラベル
\end{figure}



\subsection{中継機の構成図}
中継機デバイスの構成図は，図\ref{kan}の通りである．指揮デバイスから信号を受け取った際に，中継機である観覧車デバイスが作動する．指揮デバイスから送信されたBPMに沿って観覧車は回転し，LEDMatrixにはそのBPM数が表示される仕組みである．
\begin{figure}[h]
\centering
 \centering
\includegraphics[width=10cm]{./figs/kan.png}
\caption{中継機の構成図}
\label{kan}
\end{figure}




\subsection{楽器デバイスの構成図}
最後に，楽器デバイス受光機構のシステム構成図を図\ref{fig:jusin_sys}に示す．
本機構は，回転している観覧車デバイスのレーザー光を検知することで楽器の同期を維持する．
あらかじめ設定した位置に各楽器デバイスを設置し，観覧車デバイスのレーザーが通過した瞬間に光センサがこの光を検知することで，
演奏の同期およびBPMの変更が行われる．
% また，観覧車デバイスの中心にある中継機Arduinoからの無線通信により音量情報を受信することで楽曲の音量が変更できる．
\begin{figure}[h] % h:ここ, t:上, b:下, p:独立ページ
        \centering % 中央揃え
        \includegraphics[width=0.8\columnwidth]{figs/jusin_sys3.pdf} % 横幅を列幅の80%に指定
        \caption{楽器デバイスの受光機構のシステム構成図} % 図の説明
        \label{fig:jusin_sys} % 本文中で参照するためのラベル
\end{figure}




%%%%%%%%%%%%%%%%%%%%%%%%
\section{操作インターフェース設計}
%画面やスイッチなどのインターフェースの設計図を記載する．自分が操作したらどうなるかの説明
本節では指揮デバイスの設計および指揮デバイスの演奏開始トリガー送信機構とBPM変更機構および音量変更機構の操作インターフェース設計，中継機デバイスの表示設計について説明する．なお，楽器デバイスのの受光部はユーザの直接操作が存在しないため，操作インターフェース設計の対象外とする．

\subsection{指揮デバイスの設計}
指揮デバイスの設計について図\ref{fig:inta_siki_mitame}に示す．
本デバイスは，実際の指揮棒のようにユーザが片手で振って操作できるよう，細長い筒状の外装を採用する．
また，内部のブレッドボードに加速度センサ，可変抵抗器，ボタンを設置し，ユーザが操作しやすいようスライダー部分は外部へ露出させる
構造とする．さらに，指揮デバイスの先端部に加速度センサを，後端部に外部電源を設置することで，操作時の重心を安定させユーザビリティを向上させる．

\begin{figure}[h] % h:ここ, t:上, b:下, p:独立ページ
        \centering % 中央揃え
        \includegraphics[width=0.8\columnwidth]{figs/siki_mitame2.pdf} % 横幅を列幅の80%に指定
        \caption{指揮デバイス設計} % 図の説明
        \label{fig:inta_siki_mitame} % 本文中で参照するためのラベル
\end{figure}


\subsection{指揮デバイスの演奏開始トリガー送信機構の設計}
次に，指揮デバイスの演奏開始トリガー送信機構の設計について図\ref{fig:inta_kaisi}に示す．
本機構は，ユーザがボタンを押下することで演奏開始トリガーを送信する．
ボタンの押下によって変化する電圧を指揮デバイスで検知し，観覧車デバイスへ送信をし観覧車デバイスが動作を開始する．
% そして，観覧車デバイスが回転を開始し，そのLED光を楽器デバイスの受光部が
% 検知することで，楽曲が再生される仕様である．
\begin{figure}[h] % h:ここ, t:上, b:下, p:独立ページ
        \centering % 中央揃え
        \includegraphics[width=0.8\columnwidth]{figs/kaisi_inta2.pdf} % 横幅を列幅の80%に指定
        \caption{演奏開始トリガー送信機構設計} % 図の説明
        \label{fig:inta_kaisi} % 本文中で参照するためのラベル
\end{figure}


\subsection{指揮デバイスのBPM変更機構および音量変更機構の設計}
次に，指揮デバイスのBPM変更機構および音量変更機構の設計について図\ref{fig:inta_siki},図\ref{fig:inta_siki2}に示す．本機構は，ユーザがスライド式可変抵抗器を操作することで楽曲のBPMおよび音量が変更できる．
ユーザがスライド式可変抵抗器をスライドし，指揮棒を振ると新しい設定BPM情報が観覧車デバイスへ無線送信される．これにより，観覧車デバイスの回転速度がこの値に追従して変化し，その
回転を介して各楽器デバイスへ情報が伝達されることで，楽器のBPMが変更される．この一連の処理によって観覧車デバイスの回転速度と楽曲のBPMが常に同期される仕様である．
また音量変更は，BPMと同様にユーザがスライド式可変抵抗器をスライドすると各楽器用Arduinoに情報が送信され変更される．また，楽曲のBPMおよび音量は五段階に分けて変更できる．


\begin{figure}[htbp] % h:ここ, t:上, b:下, p:独立ページ
        \centering % 中央揃え
        \includegraphics[width=0.8\columnwidth]{figs/siki_inta2.pdf} % 横幅を列幅の80%に指定
        \caption{BPM変更機構設計} % 図の説明
        \label{fig:inta_siki} % 本文中で参照するためのラベル
\end{figure}

\begin{figure}[htbp] % h:ここ, t:上, b:下, p:独立ページ
        \centering % 中央揃え
        \includegraphics[width=0.8\columnwidth]{figs/siki_inta3.pdf} % 横幅を列幅の80%に指定
        \caption{音量変更機構設計} % 図の説明
        \label{fig:inta_siki2} % 本文中で参照するためのラベル
\end{figure}


\subsection{中継機デバイスの表示}
最後に中継機デバイスの表示設計について説明する．中継機デバイスでは，受信したBPM情報に従って観覧車を回転させる．
また，LEDMatrixを用いて現在のBPM値を表示することで，
ユーザがテンポを視覚的に確認できるようにする．
動作に関しては，上記の図\ref{kan}の通りである．


%%%%%%%%%%%%%%%%%%%%%%%%
\section{状態遷移図}
\subsection{全体の状態遷移図}
今回提案するシステムの状態遷移図を説明する．状態遷移図は図\ref{jyoutai}の通りである．
電源を入れ，ボタンを押下すると観覧車が作動する状態になる．その後回転する際に観覧車にあるレーザーから光を受光し，メロディが順に流れていき演奏する状態になる．その間にBPM変更を行った場合，もう一度ボタンを押下するとBPMを更新し，回転速度が変化する．演奏終了すると待機状態に戻り，そのまま電源が切れる流れである．
\begin{figure}[h]
\centering
 \centering
\includegraphics[width=13cm]{./figs/jyoutai.png}
\caption{状態遷移図}
\label{jyoutai}
\end{figure}


\subsection{指揮デバイスの状態遷移図}
次に，指揮デバイスにおける状態遷移について説明する．
まず，指揮デバイスのシーケンス図を図\ref{fig:jou_siki}に示す．
ユーザが指揮デバイスを上下に振ると加速度センサが動作を検知しArduinoに送信される．この時，値が閾値を超過した場合，
観覧車デバイスと無線通信を開始し観覧車デバイスが送信した演奏情報を基に回転する．また，ボタンを押下した時，各楽器デバイスへ楽曲の初期BPMおよび音量情報が送信されることで演奏が開始される．

\begin{figure}[htbp] % h:ここ, t:上, b:下, p:独立ページ
        \centering % 中央揃え
        \includegraphics[width=0.8\columnwidth]{figs/siki_jou5.pdf} % 横幅を列幅の80%に指定
        \caption{指揮デバイスのシーケンス図} % 図の説明
        \label{fig:jou_siki} % 本文中で参照するためのラベル
\end{figure}

\subsection{楽器デバイス受光機構の状態遷移図}
次に，楽器デバイスの受光機構のシーケンス図を図\ref{fig:jou_jusin}に示す．
観覧車デバイスのレーザー光が初回に通過した際，光センサが楽器デバイス内のArduinoへデジタル信号を送信し，PCから楽器の音色が再生される．
演奏中は観覧車デバイスのレーザー光が楽器デバイスの受光部を通過するごとに光センサがArduinoへデジタル信号を送信する．この通過間隔から
BPMを算出することで，最新のBPM値をPCへ送信する．

\begin{figure}[htbp] % h:ここ, t:上, b:下, p:独立ページ
        \centering % 中央揃え
        \includegraphics[width=0.8\columnwidth]{figs/jusin_jou5.pdf} % 横幅を列幅の80%に指定
        \caption{楽器デバイスの受光機構のシーケンス図} % 図の説明
        \label{fig:jou_jusin} % 本文中で参照するためのラベル
\end{figure}






%%%%%%%%%%%%%%%%%%%%%%%%
%%%%%%%%%%%%%%%%%%%%%%%%
%%%%%%%%%%%%%%%%%%%%%%%%
\chapter{システムの詳細設計}
%%%%%%%%%%%%%%%%%%%%%%%%
\section{部品一覧}
本節では，指揮デバイスおよび中継機デバイス，楽器デバイスの受光部を構成する部品について説明する．
指揮デバイスを構成する部品について表\ref{table:siki}に示す．
デジタル3軸加速度センサは指揮デバイスの動きの検知に使用し，スライド式可変抵抗器は
BPMおよび音量変更に使用される．また，指揮デバイスの電力供給源として$\SI{9}{V}$の乾電池を外部電源として使用する．
\begin{table}[h]
\caption{指揮デバイスに使用な機器}
\centering
\label{table:siki}
\begin{tabular}{|l|l|r|}\hline
機材 & 品番 & 個数 [個] \\ \hline
Arduino UNO R4 WiFi& - & 1 \\
スライド式可変抵抗器 10 kΩ &-& 2 \\
デジタル3軸加速度センサ & GY-291 & 1 \\
	タクトスイッチ & DTS-63 & 1 \\
	抵抗器　10k\Omega & - & 1 \\
外部電源 9 V &-& 1 \\
ジャンパーワイヤー & - & 適宜 \\ \hline
\end{tabular}
\end{table}

中継機デバイスで使用する部分を表\ref{kizai}にまとめる．この内，個別で購入する必要があるものはハイパワーギヤーボックス，ブレッドボード用DCジャックDIP化キット，ボタン電池，ボタン電池ホルダー，QIケーブル，有極コンデンサーである．
\begin{table}[h]
\caption{中継機デバイスに必要な機材}
\centering
\label{kizai}
\begin{tabular}{|l|l|r|}\hline
機材&品番&個数[個]\\ \hline
Arduino Uno R4 WiFi&-& 1\\
ルーター& -& 1\\ 
ブレッドボード & -&1\\
赤色ドットレーザーモジュール & FU650AD5-C6&8\\ 
LR44 電池ボックス リード線つき & - & 8 \\
ボタン電池 & LR44&16\\ 
9Vアルカリ電池 &-& 1\\ 
モータードライバーモジュール &AE-DRV8835-S& 1\\ 
ハイパワーギヤーボックス HE &72003& 1\\ 
ジャンパーワイヤー &-& 適宜\\ 
ブレッドボード用2.1mm標準DCジャックDIP化キット&AE-DC-POWER-JACK-DIP & 1 \\
超小型スイッチングACアダプター 5 V 1 A &-& 1\\
QIケーブル2P-1Px2 &311-262&1\\
抵抗器 10 Ω &- &1\\ 
抵抗器 100 Ω &-& 8\\
コンデンサー 0.01 \mu F&-&  1\\ 
コンデンサー 0.047 \mu F& -& 1\\ 
有極性コンデンサー 100 \mu F &-&  1\\ \hline
\end{tabular}
\end{table}

次に，楽器デバイスが受光するための部品について表\ref{table:ota}に示す．
光センサモジュールとしてCdSセルを使用する．
CdSセルにレーザー光が照射された際の出力をArduinoのアナログピンへ接続することで，観覧車デバイスのレーザー光の通過を検知する．また，今回はドラム，ピアノ，ストリングス，フルートの4種類の
楽器の音色を再現するためにそれぞれにArduinoを使用するため，光センサおよびブレッドボードも4台用意する．
\begin{table}[h]
\caption{楽器デバイスに使用な機器}
\centering
\label{table:ota}
\begin{tabular}{|l|l|r|}\hline
機材 & 品番 & 個数 [個] \\ \hline
Arduino UNO R4 & - & 4 \\
CdSセル  &GL5516& 4 \\
ブレッドボード & -& 4 \\
ジャンパーワイヤー & - & 適宜 \\ \hline
\end{tabular}
\end{table}


\clearpage
%%%%%%%%%%%%%%%%%%%%%%%%
%%%%%%%%%%%%%%%%%%%%%%%
\section{ハードウェア設計}
\subsection{指揮デバイス}
指揮デバイスの回路図を図\ref{fig:kairo_siki} に示す．
電源仕様は，$\SI{9}{V}$乾電池をArduinoのVINピンに入力することで，Arduino 内蔵の電圧レギュレータにより$\SI{5}{V}$および$\SI{3.3}{V}$ に降圧し，各モジュールへ安定した電力を供給する．

加速度センサにはGY-291（ADXL345）を用いており，I2C 通信により Arduino と接続する．CS ピンを Arduino の 3.3 V に接続することで I2C モードに設定し，さらに SDO ピンを GND に接続することで I2C アドレスを 0x53 に固定している．また，SDA ピンおよび SCL ピンはそれぞれデータ通信線およびクロック信号線として Arduino の A4，A5 ピンに接続している．加えて，加速度が内部設定した閾値を超過した場合には，INT1 ピンから Arduino の D2 ピンへ割り込み信号が出力される．これにより，指揮デバイスの振り動作を検知し，モータ制御用 Arduino との UDP 通信を開始する．

さらに，演奏開始および終了信号の送信に用いるボタン入力判定回路を実装する．ボタンの誤判定を防止するため，10 kΩ のプルアップ抵抗を用いた回路を構成する．ボタンが押されていない状態では，入力ピン D3 はプルアップ抵抗を介して 5 V に接続されるため HIGH 状態となる．一方，ボタン押下時には回路が GND に短絡され，入力電位が LOW に遷移する．指揮デバイスのプログラムでは D3 ピンの状態変化を常時監視することで，演奏開始および終了操作を判定する．

\begin{figure}[htbp] % h:ここ, t:上, b:下, p:独立ページ
        \centering % 中央揃え
        \includegraphics[width=0.8\columnwidth]{figs/kairo_siki.png} % 横幅を列幅の80%に指定
        \caption{指揮デバイスの回路図} % 図の説明
        \label{fig:kairo_siki} % 本文中で参照するためのラベル
\end{figure}
\subsection{観覧車型中継機デバイス}
中継機の詳細設計について説明する．図\ref{metro}について，モータードライバーにArduino D5，D6ピンを接続し，PWM制御を行えるようにする．モータードライバーへの電力供給はArduinoの5Vからと，DCジャックからのACアダプター5 V 1 Aを用いる．モータードライバー及びDCジャックはピンヘッダーのはんだ付けが必要なため，はんだ付けを行う．ギアボックスは0.01 \mu Fのコンデンサーと直接はんだ付けを行う．

次に，観覧車の設計について説明する．作成するために，3Dプリンタを用いて作成する．
観覧車の直径は約20 cmとし，ボタン電池付きレーザーを45度間隔で配置する．そして，観覧車を支える左右の支柱および土台を設計する．
観覧車本体は，
以下のパーツに分割して設計を行う．
\begin{itemize}

	\item 土台 \\
	観覧車全体を支える部分であり，
	内部にArduino Uno R4 WiFi，
	モータードライバー，
	ギヤボックスモータ，
	DCジャックおよび配線を収納する．
	また，支柱を固定する役割も持つ．

	\item 支柱 \\
	観覧車の回転軸を左右から支える部品である．
	中央ハブおよび回転軸を安定して保持し，
	観覧車が滑らかに回転できるようにする．

	\item 回転軸 \\
	ギヤボックスモータの回転を
	観覧車本体へ伝達する部品である．
	モータと中央ハブを接続し，
	観覧車全体を回転させる役割を持つ．

	\item 中央ハブ \\
	観覧車中心部の接続部品であり，
	回転軸と回転リングを固定する役割を持つ．
	回転力を観覧車全体へ伝える重要な部品である．

	\item 回転リング \\
	観覧車の円形部分であり，
	ゴンドラを取り付ける土台となる．
	モータによって回転し，
	レーザーによる光信号送信を行う部分でもある．
\end{itemize}

\begin{figure}[h]
\centering
 \centering
\includegraphics[width=17cm]{./figs/kairo.png}
\caption{観覧車型中継機デバイス回路図}
\label{metro}
\end{figure}


\subsection{レーザー光送信装置}
レーザー光送信装置の詳細設計について説明する．図\ref{LED}について，回路図の通りにボタン電池と抵抗，レーザーをはんだ付けする．これらを観覧車のリングに45度間隔に装着する．

\begin{figure}[h]
\centering
 \centering
\includegraphics[width=5cm]{./figs/LED.png}
\caption{LED回路図}
\label{LED}
\end{figure}


\subsection{楽器デバイス受光部装置}
楽器デバイスの受光部の回路図を図\ref{fig:kairo_hikari}に示す．
電源仕様は，PCとのUSB接続から電力を供給する．光センサはArduinoの$\SI{5}{V}$ピンから給電されるため，外部電源は不要である．
また，光センサが光を受光した瞬間にA0のアナログピンにセンサ値を出力する．
このセンサ値の変化を
Arduinoで処理されることで光を検知する仕様である．


\begin{figure}[htbp] % h:ここ, t:上, b:下, p:独立ページ
        \centering % 中央揃え
        \includegraphics[width=0.8\columnwidth]{figs/kairo_hikari.pdf} % 横幅を列幅の80%に指定
        \caption{楽器デバイスの受光部の回路図} % 図の説明
        \label{fig:kairo_hikari} % 本文中で参照するためのラベル
\end{figure}


\section{Processing}
本節では，音色の波形生成と発音を担うソフトウェア（Processing）の詳細設計について述べる．



\subsection{システムアーキテクチャとファイル構成}

Processingのプログラムは可読性を考慮し，機能ごとにファイルを分割している．ファイル構成を表\ref{tab:ProcessingFile}に示す．
\begin{table}[H]
    \caption{Processingのファイル構成}
    \label{tab:ProcessingFile}
    \centering
    \begin{tabular}{|l|p{10cm}|}
        \hline
        \textbf{ファイル名} & \textbf{説明} \\ \hline
        MainCode.pde & プログラムを総括する\texttt{SystemController}． \\ \hline
        DataUtils.pde & 波形生成に使用するJSONファイルのデータを扱う\texttt{DataUtils}． \\ \hline
        InstrumentPlayer.pde & 発音命令を委譲し出力する\texttt{InstrumentPlayer}． \\ \hline
        (楽器名).pde & 各楽器の波形を生成する． \\ \hline
        data / xxx.json & JSONファイルはdataディレクトリに格納する． \\ \hline
    \end{tabular}
\end{table}
MainCode.pdeでは，Processingのプログラム全体を統括し，\texttt{setup}メソッドや\texttt{serialEvent}メソッドなどを記載する．
DataUtils.pdeでは，dataディレクトリに格納されたJSONファイルを読み込み，配列として保持する．
この配列は，\texttt{GetDynamicProfile}メソッドや，文字列でデータ種別を指定する統合ゲッター\texttt{GetLUT}メソッドを介してのみ取得できる設計とする．
InstrumentPlayer.pdeでは，発音命令を受けて(楽器名).pdeにある各楽器の波形生成用クラスを呼び出し，生成した波形を出力する．

本システムは複数のクラス群によって構成されており，その全容を図\ref{fig:processingAllClass}に示す．
\begin{figure}[H]
  \centering
  \includegraphics[width=1\linewidth]{figs/2026HCK_processing_クラス図-ページ-1.drawio.png}
  \caption{Processingのクラス図}
  \label{fig:processingAllClass}
\end{figure}
全体の制御を担う\texttt{SystemController}がシリアル通信の受信や描画処理を統括し，
発音命令は\texttt{InstrumentPlayer}クラスを経由して各楽器固有のクラス（\texttt{StringsInstrument}，\texttt{FluteInstrument} など）へ委譲される．
また，合成に必要な倍音データなどは \texttt{DataUtils} クラスに集約されている．



\subsection{通信プロトコルとエラーハンドリング}
ArduinoからPCへのデータ送信には有線のシリアル通信を採用するものとし，
その形式と送信タイミングを表\ref{tab:headerN}，表\ref{tab:headerB}，表\ref{tab:headerV}に示す．
ここで，本提案では動的なテンポ変化に対応するため，楽曲の進行を測る指標として，msなどの実時間ではなくTickを用いる．Tickとは時間とBPMによって算出される値であり，BPMによって音符の実際の長さが変化しても，相対的に楽曲を管理できる\cite{bib:tick}．

% --- 発音指示パケット ---
\begin{table}[H]
    \caption{発音指示パケット（ヘッダ N）}
    \label{tab:headerN}
    \centering
    \begin{tabular}{|l|p{10cm}|}
        \hline
        \textbf{形式} & \texttt{N,pitch1,duration1,startVel1,endVel1,\newline pitch2,duration2,startVel2,endVel2,... \textbackslash n} \\ \hline
        \textbf{例} & \texttt{N,60,48,80,112,64,48,80,112\textbackslash n} \newline （ドとミの和音を長さ48Tickで，共にmfからffへ変化させながら同時に鳴らす） \\ \hline
        \textbf{送信タイミング} & Arduinoの現在Tickが，楽譜の\texttt{timing}に到達した瞬間． \\ \hline
    \end{tabular}
\end{table}

% --- BPM変更指示パケット ---
\begin{table}[H]
    \caption{BPM変更指示パケット（ヘッダ B）}
    \label{tab:headerB}
    \centering
    \begin{tabular}{|l|p{10cm}|}
        \hline
        \textbf{形式} & \texttt{B,bpmValue\textbackslash n} \\ \hline
        \textbf{例} & \texttt{B,130\textbackslash n} \\ \hline
        \textbf{送信タイミング} & 指揮からBPMの変更を受信した瞬間． \\ \hline
    \end{tabular}
\end{table}

% --- ボリューム変更指示パケット ---
\begin{table}[H]
    \caption{ボリューム変更指示パケット（ヘッダ V）}
    \label{tab:headerV}
    \centering
    \begin{tabular}{|l|p{10cm}|}
        \hline
        \textbf{形式} & \texttt{V,volumeValue\textbackslash n} \\ \hline
        \textbf{例} & \texttt{V,99\textbackslash n} \\ \hline
        \textbf{送信タイミング} & 指揮から全体のマスターボリューム変更を受信した瞬間． \\ \hline
    \end{tabular}
\end{table}

先頭にデータの種類を表すヘッダ文字を付与したカンマ区切りの文字列とし，
ヘッダ文字がNの場合は発音指示，Bの場合はBPM変更指示，Vの場合はボリューム変更指示であると判別する．
%BPM変化に素早く対応するため，BPM変更指示は発音指示より優先的に送信する．
%また，和音の同時発音性を保証するため，同一タイミングの音符は1行のパケットにまとめて送信する．

\begin{table}[H]
    \centering
    \caption{エラーおよび境界値の処理仕様}
    \label{tab:error_handling}
    \begin{tabularx}{\textwidth}{|l|X|l|}
        \hline
        \textbf{対象項目} & \textbf{判定条件} & \textbf{処理内容} \\ \hline
        パケット (N) & 要素数からヘッダを除いた値が4の倍数でない & 破棄（欠損とみなす） \\ \hline
        パケット (B, V) & ヘッダを含む要素数が2でない & 破棄（欠損とみなす） \\ \hline
        各パラメータ & 定められた許容範囲（0～127等）を逸脱 & \texttt{constrain()} による補正 \\ \hline
    \end{tabularx}
\end{table}

本システムでは，通信環境によるパケットの欠損や想定外の数値入力に対処するため，表\ref{tab:error_handling}に示す通りのエラーハンドリングを実装する．
シリアル通信においては，ノイズ等の影響によりデータの一部が欠損する可能性がある．
特に発音指示パケット（ヘッダN）は，ノート数に応じてデータ長が可変となるため，
分割後の要素数を確認することで，受信したデータがプロトコルに準拠しているかを判定している．
不整合が検知された場合は，誤作動を防止するために当該パケットを即座に破棄する設計とした．

具体的な処理手順として，シリアル通信で受信した文字列は，
末尾に含まれる改行コード（\texttt{\textbackslash n} や \texttt{\textbackslash r}）が後のデータ変換処理において数値変換エラーを引き起こす原因となる．
そのため，パケットを受信した直後に\texttt{trim}関数を用いて空白文字や改行コードを確実に消去した上で，
\texttt{split}関数によるカンマ分割を行う仕様としている．
分割後の要素数からヘッダを除いた値が規定数（発音指示であれば4の倍数）に合致するかを確認することで，不完全なパケットを安全に弾くことができる．
BPMやボリュームなどの数値データについては，指揮側からの入力値がソフトウェア上の許容範囲を超過する場合が考えられる．
これに対し，単にエラーとして捨てるのではなく，\texttt{constrain}関数を用いて上限・下限値に丸め込みを行うことで，システムの動作継続性を優先させる．



\subsection{システム統括とデータ構造}

\subsubsection{SystemController（システム統括・状態管理）}
SystemControllerのフローチャートを図\ref{fig:flowchartSystemController}に示す．
\begin{figure}[H]
  \centering
  \includegraphics[width=1\linewidth]{figs/SystemController.pdf}
  \caption{SystemControllerのフローチャート}
  \label{fig:flowchartSystemController}
\end{figure}
システムの根幹となるメインクラスであり，Arduinoからのシリアル通信（\texttt{myPort}）の受信，
Minimライブラリを用いた音響処理の初期化，およびFFT（高速フーリエ変換）による波形解析の管理を行う．
本クラスで定義されている主要な属性（変数）を表\ref{table:system_variables}に，メソッド（関数）を表\ref{table:system_methods}に示す．
\begin{table}[H]
    \centering
    \caption{SystemControllerの主要な属性（変数）}
    \label{table:system_variables}
    \begin{tabular}{|l|l|p{7cm}|}
        \hline
        \textbf{変数名} & \textbf{型} & \textbf{説明} \\ \hline
        \texttt{minim} & \texttt{Minim} & 音響合成ライブラリMinimのメインインスタンス． \\ \hline
        \texttt{out} & \texttt{AudioOutput} & Minimの最終的な音声出力先． \\ \hline
        \texttt{fft} & \texttt{FFT} & 出力波形を周波数スペクトルに変換する解析オブジェクト． \\ \hline
        \texttt{player} & \texttt{InstrumentPlayer} & 発音命令を各楽器クラスへ委譲するための仲介インスタンス． \\ \hline
        \texttt{utils} & \texttt{DataUtils} & 波形生成用のLUTデータを保持・提供するユーティリティ． \\ \hline
        \texttt{arduinoNumber} & \texttt{byte} & 接続するArduinoの番号．演奏する楽器の識別に用いる． \\ \hline
        \texttt{musicBPM} & \texttt{short} & 現在の楽曲テンポ．Tickから秒数への変換基準となる． \\ \hline
        \texttt{currentMasterVolume} & \texttt{byte} & 現在の全体音量．最終出力時の振幅に乗算される． \\ \hline
        \texttt{isDebugMode} & \texttt{boolean} & 単体テスト用フラグ．true時にPCキーボード入力を受け付ける． \\ \hline
        \texttt{soundOctave} 等 & \texttt{byte} & デバッグモード時における手動発音用のパラメータ群． \\ \hline
    \end{tabular}
\end{table}
\begin{table}[H]
    \centering
    \caption{SystemControllerの主要なメソッド}
    \label{table:system_methods}
    \begin{tabular}{|l|l|p{2.8cm}|p{6.8cm}|}
        \hline
        \textbf{メソッド名} & \textbf{戻り値} & \textbf{引数} & \textbf{説明} \\ \hline
        \texttt{setup} & \texttt{void} & なし & 画面サイズ，Minim，通信，各インスタンスの初期化を行う． \\ \hline
        \texttt{draw} & \texttt{void} & なし & オシロスコープおよびスペクトラムアナライザをリアルタイム描画する． \\ \hline
        \texttt{serialEvent} & \texttt{void} & \texttt{Serial port} & Arduinoからのパケットを受信し，発音やBPM変更処理へ分岐する． \\ \hline
        \texttt{Play} & \texttt{void} & なし & デバッグモード時に，選択中の楽器のテスト発音を実行する． \\ \hline
        \texttt{keyPressed} & \texttt{void} & なし & デバッグモード時のPCキーボード入力を検知しパラメータを変更する． \\ \hline
    \end{tabular}
\end{table}

\subsubsection{DataUtils（データ構造とルックアップテーブル）}
DataUtilsのフローチャートを図\ref{fig:flowchartDataUtils}に示す．
\begin{figure}[H]
  \centering
  \includegraphics[width=1\linewidth]{figs/DataUtils.pdf}
  \caption{DataUtilsのフローチャート}
  \label{fig:flowchartDataUtils}
\end{figure}
本システムにおける音色合成の基礎となるのは，Pythonを用いた実波形解析によって構築されたルックアップテーブル（LUT）である．
%楽器の音声データに対しFFT（高速フーリエ変換）解析を行い，各倍音の周波数・振幅・減衰時間といったパラメータを抽出している．
%Pythonによって解析・生成されたこれらの膨大なパラメータ群は，JSON形式の外部データファイルとしてエクスポートされる．
%本クラスは，システムの初期化処理（\texttt{setup}）においてJSONファイルを読み込み，
%各楽器および強弱に応じた多次元配列としてメモリ上に展開する役割を担う．
本クラスの主要な属性を表\ref{table:datautils_variables}に，メソッドを表\ref{table:datautils_methods}に示す．
\begin{table}[H]
    \centering
    \caption{DataUtilsの属性（カプセル化されたLUTデータ群）}
    \label{table:datautils_variables}
    \begin{tabular}{|l|l|p{8cm}|}
        \hline
        \textbf{変数名} & \textbf{型} & \textbf{説明} \\ \hline
        \texttt{fluteFFLUT} 等 & \texttt{float[][]} & フルートの強弱（ff/mf/pp）ごとの倍音プロファイル配列． \\ \hline
        \texttt{violinFFLUT} 等 & \texttt{float[][]} & 各ストリングス楽器の強弱ごとの倍音プロファイル配列． \\ \hline
        \texttt{pianoFFLUT} 等 & \texttt{float[][]} & ピアノの強弱ごとの倍音プロファイル配列． \\ \hline
        \texttt{pianoDecayLUT} 等 & \texttt{float[][]} & ピアノ専用の減衰時間，うなり速度，うなり深度の配列． \\ \hline
        \texttt{kickLUT} 等 & \texttt{float[][]} & 各ドラムの周波数・強度・減衰時間を格納した配列． \\ \hline
    \end{tabular}
\end{table}
\begin{table}[H]
    \centering
    \caption{DataUtilsのメソッド}
    \label{table:datautils_methods}
    \begin{tabular}{|l|l|p{3.5cm}|p{4cm}|}
        \hline
        \textbf{メソッド名} & \textbf{戻り値} & \textbf{引数} & \textbf{説明} \\ \hline
        \texttt{DataUtils} & コンストラクタ & なし & dataディレクトリから全JSONファイルを読み込みメモリに展開する． \\ \hline
        \texttt{LoadLUT} & \texttt{float[][]} & \texttt{String filename} & JSONファイルをパースし，2次元配列に変換する内部関数． \\ \hline
        \texttt{GetProfileForPitch} & \texttt{float[]} & \texttt{float pitch, \newline float[][] lut, ...} & 指定したピッチに最も近いLUTデータを線形補間して安全に取り出す． \\ \hline
        \texttt{GetDynamicProfile} & \texttt{float[]} & \texttt{float pitch, \newline float velocity, \newline String type, ...} & ピッチと強弱の2軸から最終的な倍音構成レシピを生成・補間する． \\ \hline
        \texttt{GetLUT} & \texttt{float[][]} & \texttt{String type} & 文字列指定により目的の配列を安全に返す． \\ \hline
    \end{tabular}
\end{table}

本設計では，単一の倍音構成を使い回すのではなく，
楽器ごとに演奏の強弱（ff: フォルテッシモ，mf: メゾフォルテ，pp: ピアニッシモ）および音高（ピッチ）の二次元配列を保持する．
発音時，\texttt{GetDynamicProfile}メソッドは引数の\texttt{midiNote}（音高）と\texttt{velocity}（強弱）に基づき，これらのLUT間を線形補間演算する．
これにより，強く弾いた際の高次倍音の増加や，音域による倍音構造の変化など，楽器特有の非線形な音色変化を再現する．

また，ドラム等の打楽器群についても，周波数ピークごとの減衰時間を解析した専用のLUTを定義し，波形合成の基礎データとして利用する．
外部クラスからこれらのデータへアクセスする際は，引数の文字列によって条件分岐し，目的の多次元配列を返す\texttt{GetLUT}メソッドを実装している．
オブジェクト指向におけるカプセル化を強固にしつつ，将来的な音色や新規LUTデータの追加に対して高い拡張性を担保している．



\subsubsection{InstrumentPlayer（発音命令の委譲とスケジューリング）}
InstrumentPlayerのフローチャートを図\ref{fig:flowchartInstrumentPlayer}に示す．
\begin{figure}[H]
  \centering
  \includegraphics[width=0.5\linewidth]{figs/InstrumentPlayer.pdf}
  \caption{InstrumentPlayerのフローチャート}
  \label{fig:flowchartInstrumentPlayer}
\end{figure}
\texttt{SystemController}と各楽器の波形生成クラス（\texttt{Instrument} 系クラス）を仲介し，発音のスケジューリングを担う．
本クラスの属性およびメソッドを表\ref{table:player_variables}，表\ref{table:player_methods}に示す．
\begin{table}[H]
    \centering
    \caption{InstrumentPlayerの属性}
    \label{table:player_variables}
    \begin{tabular}{|l|l|p{8.5cm}|}
        \hline
        \textbf{変数名} & \textbf{型} & \textbf{説明} \\ \hline
        \texttt{out} & \texttt{AudioOutput} & SystemControllerから受け取ったMinimの音声出力オブジェクト． \\ \hline
    \end{tabular}
\end{table}
\begin{table}[H]
    \centering
    \caption{InstrumentPlayerのメソッド}
    \label{table:player_methods}
    \begin{tabular}{|l|l|p{4cm}|p{5cm}|}
        \hline
        \textbf{メソッド名} & \textbf{戻り値} & \textbf{引数} & \textbf{説明} \\ \hline
        \texttt{PlayFlute} & \texttt{void} & \texttt{float midiNote, \newline float duration, ...} & フルートのインスタンスを生成し\texttt{out.playNote}でスケジュールする． \\ \hline
        \texttt{PlayStrings} & \texttt{void} & \texttt{float midiNote, ...} & ストリングスのインスタンスを生成し発音をスケジュールする． \\ \hline
        \texttt{PlayPiano} & \texttt{void} & \texttt{float midiNote, ...} & ピアノのインスタンスを生成し発音をスケジュールする． \\ \hline
        \texttt{PlayKick} 等 & \texttt{void} & \texttt{float masterVol, \newline float velocity} & ドラム群のインスタンスを生成する．長さは内部エンベロープに依存するためダミー値を用いる． \\ \hline
    \end{tabular}
\end{table}

全体の出力を管理する\texttt{SystemController}からMinimの音声出力オブジェクト（\texttt{AudioOutput out}）を受け取り，コンポジションとして内部で保持する．

シリアル通信から解析されたパラメータ（音高，長さ，ベロシティなど）を受け取ると，
本クラスに定義された\texttt{PlayFlute}や\texttt{PlayPiano}といった専用メソッドが呼び出される．
メソッド内部では，指定されたパラメータを持つ楽器クラスのインスタンスを動的に生成し，
\texttt{out.playNote}を通じてMinimの内部スケジューラに登録・発音させる設計とする．

特筆すべき仕様として，管弦楽器やピアノが受信した音符の長さ（\texttt{duration}）を直接\texttt{playNote}の発音時間として渡すのに対し，
キックやスネア等の打楽器群（\texttt{PlayKick}等）については，発音の長さが楽器本来の物理的減衰に完全に依存する．
そのため，打楽器の呼び出しにおいては第2引数の\texttt{duration}にシステム上のダミー値（1.0等）を与え，
実際の消音タイミング等の制御は各楽器クラス内部のエンベロープ（\texttt{ADSR}）へ完全に一任することで，不要な処理の競合を防ぐ．



\subsubsection{パラメータ変換仕様とシステム上の制約}

Arduinoから受信したデータを，Processing内で音響処理を行うための物理量へと変換する共通ロジックについて述べる．

1小節を96Tickとした場合，四分音符は24Tickとなる．
受信した長さ（\texttt{duration}）は，現在の\texttt{musicBPM}を用いて式(\ref{eq:tick_to_sec})により実際の発音時間（秒）へと変換される．
\begin{equation}
    発音時間(sec) = duration \times \frac{60}{musicBPM \times 24} \label{eq:tick_to_sec}
\end{equation}
また，受信した0〜127の強さ（\texttt{startVelocity}および\texttt{endVelocity}）は，\texttt{map}関数を用いて0.0から1.0の振幅値へと変換される．
この値は各楽器の音量制御だけでなく，前述の\texttt{DataUtils}を介した音色決定の演算基準としても機能する．
最終的な出力波形に対しては，\texttt{masterVolume}を99.0で除算した係数が乗算され，指揮機からの全体音量指示が適用される．

各楽器クラスの\texttt{noteOn}が呼び出された瞬間に発音時間（秒）がスケジュールされる．
そのため，長い音符の発音途中でBPM変更指示（Bパケット）を受信した場合，
既に発音中の音符の終了タイミングは「変更前のBPM基準」で維持されるというシステム上の制約が存在する．
これを回避して完全なリアルタイム同期を行うには，Arduinoからの送信を「発音開始」と「終了」の2パケットに分割し，Processing側で状態管理を行う必要がある．
しかし，この手法は通信量の倍増による遅延リスクや，Minimの内部スケジューラとの競合を引き起こす可能性が高い．
実運用上，BPM変化後における「次の音符」の発音タイミングはArduinoの最新Tickによって正確に制御されるため，楽曲全体の同期は破綻しない．
以上の理由から，本システムでは動作の安定性を優先し，発音開始時に一括で秒数変換を行う方式を基本仕様とする．





\subsection{各楽器の実装仕様（Instrument系クラス）}

生成されたパラメータは\texttt{InstrumentPlayer}クラスに渡され，指定された楽器のコンストラクタが呼び出される．
全楽器クラスは，発音処理を開始する\texttt{noteOn}と，消音処理を担う\texttt{noteOff}を共通して実装する設計とする．

現実の楽器の音色に特有の生々しさや張りを再現するため，本システムの加算合成においては，完全な整数倍音ではなく，
物理的な弦の太さや管の剛性に起因する「インハーモニシティ（非整数倍音）」の数理モデルを全楽器に適用している\cite{インハーモニシティ式}．
第$n$次倍音の実際の周波数$f_n$は，式(\ref{eq:インハーモニシティ})により算出される．
\begin{equation}
    f_n = f_0 \times n \times \sqrt{1 + B \times n^2} \label{eq:インハーモニシティ}
\end{equation}
\begin{itemize}
    \item $f_n$: 第$n$次倍音の実際の周波数
    \item $n$: 倍音の次数（$n=1$は基音）
    \item $f_0$: 弦の硬さを考慮しない場合の理想的な基音周波数
    \item $B$: インハーモニシティ係数
\end{itemize}
この係数$B$は，各楽器の物理特性に合わせて個別に設定され，オシレーターの初期位相をランダム化することによるピーク割れの防止処理と共に，合成波形の基盤として組み込む．


\subsubsection{フルートの実装（FluteInstrument）}
FluteInstrumentのフローチャートを図\ref{fig:flowchartFluteInstrument}に示す．
\begin{figure}[H]
  \centering
  \includegraphics[width=1\linewidth]{figs/FluteInstrument.pdf}
  \caption{FluteInstrumentのフローチャート}
  \label{fig:flowchartFluteInstrument}
\end{figure}
フルートクラスにおける主要な属性を表\ref{table:flute_variables}に，メソッドを表\ref{table:flute_methods}に示す．
\begin{table}[H]
    \centering
    \caption{FluteInstrumentの主要な属性}
    \label{table:flute_variables}
    \begin{tabular}{|l|l|p{9cm}|}
        \hline
        \textbf{変数名} & \textbf{型} & \textbf{説明} \\ \hline
        \texttt{numHarmonics} & \texttt{byte} & 生成する倍音の数を指定する． \\ \hline
        \texttt{masterEnv} & \texttt{ADSR} & 楽器全体の音の長さを管理するマスターエンベロープ． \\ \hline
        \texttt{harmonicEnvs} & \texttt{ADSR[]} & 各倍音独立のアタック・ディケイを管理する配列． \\ \hline
        \texttt{breathEnv} & \texttt{ADSR} & 発音直後の破裂音（アタックノイズ）の長さを管理する． \\ \hline
        \texttt{sustainEnv} & \texttt{ADSR} & 持続ノイズ（サブハーモニクス層）の長さを管理する． \\ \hline
    \end{tabular}
\end{table}
\begin{table}[H]
    \centering
    \caption{FluteInstrumentのメソッド}
    \label{table:flute_methods}
    \begin{tabular}{|l|l|p{2.5cm}|p{5cm}|}
        \hline
        \textbf{メソッド名} & \textbf{戻り値} & \textbf{引数} & \textbf{説明} \\ \hline
        \texttt{FluteInstrument} & (コンストラクタ) & \texttt{float midiNote, \newline float startVel, ...} & LUTデータ取得，インハーモニシティ算出，ノイズ生成，LFO結線などの初期化を行う． \\ \hline
        \texttt{noteOn} & \texttt{void} & \texttt{float dur} & 各エンベロープ（ADSR）をトリガーし音を出力に出力する． \\ \hline
        \texttt{noteOff} & \texttt{void} & なし & 鍵盤が離された際のリリース処理とリソース解放を行う． \\ \hline
    \end{tabular}
\end{table}

フルートには最大15倍音の加算合成を採用し，インハーモニシティ係数 $B=0.0005$ を適用する．
特有の息の揺らぎを再現するため，ランダムな周期を呼吸スピードとして定義し，極小のピッチ揺らぎ（FM変調）を与える．
さらに実測データの変動率（\%）から導き出された「倍音ごとの自然な振幅揺らぎ（AM変調）」の深度を配列として保持し，
呼吸スピードに同期させて各倍音の音量を揺らすことで，人間が演奏しているような音色を再現する．

管楽器の物理特性を模倣するノイズは，アタックと持続の2層構造とする．
発音の瞬間には，ホワイトノイズにフィルターをかけたアタックノイズを付加し，息がエッジに当たる破裂音を表現する．
また，発音中はピンクノイズにローパスフィルターを適用した持続ノイズ（サブハーモニクス層）を付加する．
この持続ノイズに対しても独立したLFOで振幅揺らぎを加えることで，管内部で持続する複雑な空気の渦を再現する．

\subsubsection{ストリングスの実装（StringsInstrument）}
StringsInstrumentのフローチャートを図\ref{fig:flowchartStringsInstrument}に示す．
\begin{figure}[H]
  \centering
  \includegraphics[width=1\linewidth]{figs/StringsInstrument.pdf}
  \caption{StringsInstrumentのフローチャート}
  \label{fig:flowchartStringsInstrument}
\end{figure}
ストリングスクラスにおける主要な属性を表\ref{table:strings_variables}に示す．メソッド構成は表\ref{table:flute_methods}と同様である．
\begin{table}[H]
    \centering
    \caption{StringsInstrumentの主要な属性}
    \label{table:strings_variables}
    \begin{tabular}{|l|l|p{9cm}|}
        \hline
        \textbf{変数名} & \textbf{型} & \textbf{説明} \\ \hline
        \texttt{numHarmonics} & \texttt{byte} & 生成する倍音の数を指定する． \\ \hline
        \texttt{masterEnv} & \texttt{ADSR} & ストリングス全体の音の長さを管理するエンベロープ． \\ \hline
        \texttt{vibDepthLines} & \texttt{Line[]} & 発音後にビブラートの深さを徐々に深くするための制御用Line． \\ \hline
        \texttt{biteEnv} & \texttt{ADSR} & 発音直後の弓の引っ掛かり（摩擦）ノイズを管理する． \\ \hline
        \texttt{hissEnv} & \texttt{ADSR} & 持続的な擦過ノイズを管理する． \\ \hline
    \end{tabular}
\end{table}

ストリングスには最大30倍音の加算合成を採用する．ストリングスの広大な音域を自然に網羅するため，
音高（MIDIノート番号）に応じて4種類の弦楽器（コントラバス，チェロ，ビオラ，ヴァイオリン）のLUTプロファイルを滑らかにクロスフェードさせる処理を実装する．
具体的には，MIDIノート36（C2）以下はコントラバス単体の音色とし，そこから音高が上がるにつれて，
チェロ（C2〜C3），ビオラ（C3〜C4），ヴァイオリン（C4〜C5）へと段階的かつ連続的にブレンド割合を変化させ，
72（C5）以上ではヴァイオリン単体となるよう各倍音成分を合成する．
これにより，全音域において単一の波形を拡張した際に生じる不自然さを排除し，実際のオーケストラのような重厚かつ自然な響きの移り変わりを再現する．

さらに，インハーモニシティ係数$B$を固定値とせず，音域に基づく楽器割合から基準値を算出した上で，音高の2乗に比例して動的に増幅させる数理モデルを組み込む．
これにより，高音域になるほど高次倍音が鋭くなるハイポジション（指板の端で弾いた際の張り詰めた緊張感）を再現する．

ビブラートについては，音高に応じて揺れの速度および深さを動的に決定し，
発音直後から音符の長さに合わせて徐々に最大深度へ到達するよう遅延処理を施す．
このビブラート周期は，高次倍音ほど激しく揺れる音量揺らぎ（AM変調）とも同期する．

加えて，弓と弦の摩擦音を2層のノイズで表現する．
発音直後の引っ掛かり音（Bite）にはローパスフィルターをかけたホワイトノイズを用い，音高が上がるにつれてカットオフ周波数を高く，ディケイ時間を短くする動的制御を行う．
持続的な擦過音（Hiss）にはハイパスフィルターをかけたピンクノイズを採用し，これらもビブラートに同期して音量を揺らすことで，弓の動きを表現する．

\subsubsection{ピアノの実装（PianoInstrument）}
PianoInstrumentのフローチャートを図\ref{fig:flowchartPianoInstrument}に示す．
\begin{figure}[H]
  \centering
  \includegraphics[width=1\linewidth]{figs/PianoInstrument.pdf}
  \caption{PianoInstrumentのフローチャート}
  \label{fig:flowchartPianoInstrument}
\end{figure}
ピアノクラスにおける主要な属性を表\ref{table:piano_variables}に示す．
\begin{table}[H]
    \centering
    \caption{PianoInstrumentの主要な属性}
    \label{table:piano_variables}
    \begin{tabular}{|l|l|p{8.5cm}|}
        \hline
        \textbf{変数名} & \textbf{型} & \textbf{説明} \\ \hline
        \texttt{numHarmonics} & \texttt{byte} & 生成する倍音の数を指定する． \\ \hline
        \texttt{masterEnv} & \texttt{ADSR} & ダンパーフェルトの役割を果たし，離鍵時に即座にミュートする． \\ \hline
        \texttt{hammerEnv} & \texttt{ADSR} & ハンマーが弦を叩く一瞬の打撃ノイズを管理する． \\ \hline
        \texttt{damps} & \texttt{ArrayList<Damp>} & 指数関数的減衰を表現するDampオブジェクトのリスト．芯の弦と共鳴弦を動的に格納する． \\ \hline
    \end{tabular}
\end{table}

ピアノには最大40倍音の加算合成を採用し，打弦楽器としてのハンマーの物理挙動と弦の共鳴現象を再現する設計とする．
メソッド構成は表\ref{table:flute_methods}と同様であるが，管楽器等とは異なり発音後に自ら音量をコントロールすることは不可能であるため，
音色の決定や音量制御には開始時の\texttt{velocity}（\texttt{startVelocity}）のみを使用し，\texttt{endVelocity}は引数としない仕様とする．

減衰の表現においては，各倍音に対して直線的なADSRではなく，自然界のエネルギー減衰を再現する\texttt{Damp}クラス（指数関数的減衰）を用いる．
さらに，芯となる弦に加えわずかに周波数をずらした共鳴弦を並列で合成することで，ピアノ特有の複雑なうなりとリバーブ成分を生成する．
これら各倍音の減衰時間やうなりのパラメータは，事前解析したLUTデータを優先して適用し，データ欠損時はピッチに基づく数式モデルで自動補完する仕組みを実装する．

また，発音時のアタック成分として，打鍵の強さと音高に連動してカットオフ周波数が変化するローパスフィルター付きのホワイトノイズを付加し，
ハンマーが弦を叩く瞬間の鋭い衝撃音を再現する．
離鍵時（\texttt{noteOff}）には，全体の出力を統括するマスターエンベロープが短いリリース時間で即座にミュート処理を行い，
ダンパーが弦を押さえる挙動を再現する．

\subsubsection{ドラムの実装（Kick，Snare，HiHatInstrument）}
ドラムのフローチャートを図\ref{fig:flowchartDrumInstrument}に示す．
\begin{figure}[H]
  \centering
  \includegraphics[width=1\linewidth]{figs/DrumInstrument.pdf}
  \caption{ドラムのフローチャート}
  \label{fig:flowchartDrumInstrument}
\end{figure}
ドラムパートの各クラスの代表的な属性を表\ref{table:drum_variables}に示す．
\begin{table}[H]
    \centering
    \caption{ドラム系クラスの主要な属性}
    \label{table:drum_variables}
    \begin{tabular}{|l|l|p{9cm}|}
        \hline
        \textbf{変数名} & \textbf{型} & \textbf{説明} \\ \hline
        \texttt{masterEnv} & \texttt{ADSR} & ドラム全体の余韻を管理するエンベロープ（HiHatでは極短に設定）． \\ \hline
        \texttt{harmonicEnvs} & \texttt{ADSR[]} & LUTから取得した周波数ごとの独立した減衰時間を管理する． \\ \hline
        \texttt{noiseEnv} & \texttt{ADSR} & スナッピーや金属の摩擦成分など，減算合成ノイズの寿命を管理する． \\ \hline
    \end{tabular}
\end{table}

ドラムパートは各楽器の物理特性に合わせ，発音時に\texttt{DataUtils}の\texttt{GetLUT}メソッドから取得した周波数・強度・減衰時間データを用いた加算合成と，ノイズによる減算合成を組み合わせた構成とする．
各周波数成分には個別の減衰時間を持つエンベロープを適用し，リアルな胴鳴りと余韻を生成する．

キック（\texttt{KickInstrument}）は，LUT成分ごとの加算合成に加え，
発音の瞬間に周波数を急激に降下させることで，ビーターが打面を叩くアタック感を生成する．
スネア（\texttt{SnareInstrument}）は，LUTの成分と7750Hz以上のハイパスフィルターを用いたノイズ成分を合成する．
ハイハット（\texttt{HiHatInstrument}）は，LUT成分と10000Hz以上のハイパスフィルタを用いたノイズ成分を合成する．
全体を統括するマスターエンベロープのリリース時間を0.05秒に設定することで，クローズハイハットの鋭くキレのある刻み音を実現する．





\subsection{デバッグおよび単体テスト機能（Processing）}

開発および音色調整フェーズにおいて，Arduinoからのシリアル通信が確立していない状態でもPC単体で検証を行えるよう，
\texttt{SystemController}内部にデバッグモードを実装する．

\texttt{isDebugMode}がtrueの場合，\texttt{keyPressed}メソッドを介してPCのキーボード入力を受け付ける．
キーボードの特定の文字キーを音階（A〜J等）や楽器切り替え（Z，X等）に割り当て，手動での\texttt{Play}メソッド呼び出しを可能とする．
また，\texttt{draw}メソッド内において\texttt{Minim}の出力波形（オシロスコープ）および\texttt{FFT}オブジェクトを用いた
周波数スペクトラム（スペクトラムアナライザ）を画面上にリアルタイム描画し，合成された倍音構成やノイズ成分の視覚的な解析を支援する設計とする．
%%%%%%%%%%%%%%%%%%%%%%%%
\section{指揮デバイスのセンサ制御}
\subsection{ボタン入力の判定}
指揮デバイスにおいて，演奏の開始および終了信号の送信に使われるボタンの入力判定処理を実装する．

ボタンの誤判定を防ぐため，プルアップ抵抗を用いた回路を構成している．ボタンが押されていない状態では，回路は電源電圧にプルアップされているため，該当するデジタルピンには常に5Vが印加される．

一方，ボタンが押下された状態では，回路がGNDに短絡されるため，ピンの電位はLOWへと遷移する．指揮デバイスのプログラムでこのデジタルピンの電圧状態の変化を常時監視・検知することにより，ボタンの押下を判定する．
\subsection{加速度センサ制御}
指揮デバイスは，加速度センサを用いてユーザが棒を振ったことを検知し，演奏情報の更新を行う
．
センサ値の急激な変化を読み取ることで動作を判定し，誤検知防止機能をつけることで，安定した動作を実現する．

加速度センサの制御で使用する変数の仕様を表\ref{tab:acc}に示す．
アルゴリズムとしては，差分算出と判定，多重検知の防止，状態の反転，事後処理の4段階に分けられる．
差分算出と判定では，現在の加速度センサ値(currentAccVal)と前回の値(lastAccVal)の差の絶対値を算出する．
ここでcurrentAccValの値は，3軸のセンサ値をそれぞれ2乗したものを加算し平方根をとった値とする．
この差分が閾値accThresholdを超え，かつ前回の検知からshakeInterval以上の時間が経過している場合に，棒を振ったと判断する．

多重検知の防止では，一度振る動作を検知すると，lastDetectTimeを現在の時刻を更新する．以降，shakeIntervalが経過するまでは，加速度の変化を無視することで，一連の動作による複数回の誤検知を防止する．

状態の反転では，振る動作を検知するたびに，isStartedのbool値を反転させる．これにより奇数回目の入力で開始(true)，偶数回目の入力で停止(false)という状態を表す．

事後処理としては，次回の判定に備えて，currentAccValの値をlastAccValに代入し，前回の処理時間を更新する．

%\begin{figure}[h]            
    %\centering                
    %\includegraphics[width=0.8\columnwidth]{./figs/sikensu.png} 
    %\caption{通信全体のシーケンス図}
    %\label{zu:sikensu}       
%\end{figure}

\begin{table}[H]
  \centering
  \caption{加速度センサの制御で使用する変数の仕様一覧}
  \label{tab:acc}
  \begin{tabular}{llp{6cm}} \hline
    変数名 & 型 & 説明 \\ \hline
	  \texttt{currentAccVal} & \texttt{int} & 加速度センサからの値を格納 \\ \hline
    \texttt{lastAccVal} & \texttt{int}  & 変化量算出のための直前の加速度センサ値を格納 \\\hline
    \texttt{accThreshold} & \texttt{const int}  & 振ったと判断するための変化量の閾値を定義 \\\hline
    \texttt{currentTime} & \texttt{unsigned long}  & millis関数を用いて，現在の時間をms単位で格納 \\ \hline
    \texttt{lastDetectTime} & \texttt{unsigned long}  & 前回振ったことを検知した時間をms単位で格納 \\ \hline
    \texttt{shakeInterval} & \texttt{const int}  & 連続検知を防止するための待機時間を定義 \\ \hline
	  \texttt{inStarted} & \texttt{bool}  & 楽器の演奏状態を保持(true:演奏中/false:停止中)\\ \hline
  \end{tabular}
\end{table}

%%%%%%%%%%%%%%%%%%%%%%%
\subsection{可変抵抗器を用いたBPM・音量読み取り機能}
指揮デバイスでは，演奏の要素であるBPMと音量を，ユーザが操作できるように可変抵抗器を用いて5段階で調整する．
アナログ入力値の特性を考慮したデータ処理を行うことで，安定した制御環境を構築する．
この処理で使用する変数の仕様を表\ref{tab:bpmvol1}とに示す．

使用するパラメータは，可変抵抗器から得られる0〜1023の値を5つの範囲に分割し，それぞれの段階を決定する．

アナログピンで取得される電圧値は，ノイズや電気的特性により細かな変動が生じやすい．
この変動が段階の境界付近で発生すると，BPMや音量が意図せず頻繁に切り替わる問題が生じる．
これを防止するために，移動平均を用いる．
複数回の電圧値を常に平均化して評価することで，電圧値の推移を平滑化し，段階決定の安定性を確保する．

アルゴリズムとしては，サンプリングと移動平均化，範囲分割，パラメータの抽出と出力の3段階に分けられる．
サンプリングと移動平均化はまず，BPMと音量変更用の可変抵抗器の電圧値をそれぞれ配列bpmSamples，volSamplesにsampleSize分だけ読み取り保存する．
配列が埋まったら，2つの配列の合計をそれぞれ変数bpmTotal，volTotalに格納し，それをsampleSizeで除算する．
この操作で算出した移動平均をそれぞれ，変数averageBpmVal，averageVolValに格納する．
また，配列がいっぱいになった後は，一番古い要素から更新していくため，変数readIndexで古い要素のインデックスを記録し，配列のそのインデックスの位置の要素を更新していく．
これらの操作により，細かな変動を抑えた電圧値を取得する．


範囲分割では，移動平均化された0〜1023の値が格納された変数averageBpmVal，averageVolValをあらかじめ表\ref{tab:range}で定義した5つの範囲に当てはめて，現在の段階を変数currentBpmStep，currentVolValに格納する．

\begin{table}[H]
  \centering
  \caption{BPM・音量読み取り機能で使用する変数の仕様一覧1}
  \label{tab:bpmvol1}
  \begin{tabular}{llp{6cm}} \hline
    変数名 & 型 & 説明 \\ \hline
	  \texttt{bpmPin} & \texttt{const int} & BPM変更用の可変抵抗器を接続するアナログ入力ピン番号を定義 \\\hline
	  \texttt{volPin} & \texttt{const int} & 音量変更用の可変抵抗器を接続するアナログ入力ピン番号を定義 \\\hline
    \texttt{sampleSize} & \texttt{const int}  & 移動平均算出のためのサンプル数を定義 \\\hline
	  \texttt{bpmSamples[]} & \texttt{int}  & BPM用のサンプル値を保持するための配列 \\\hline
	  \texttt{volSamples[]} & \texttt{int}  & 音量用のサンプル値を保持するための配列 \\\hline
	  \texttt{readIndex} & \texttt{int}  & 配列内での現在の書き込み位置を管理するインデックス \\\hline
	  \texttt{bpmTotal} & \texttt{int}  & bpmSamples[]内の累積合計を格納 \\ \hline
	  \texttt{volTotal} & \texttt{int}  & volSamples[]内の累積合計を格納 \\ \hline
    \texttt{averageBpmVal} & \texttt{int}  & 算出されたBPM用のアナログ入力の平均値を格納\\ \hline
    \texttt{averageVolVal} & \texttt{int}  & 算出された音量用のアナログ入力の平均値を格納 \\ \hline
	  \texttt{currentBpmStep} & \texttt{int}  & 算出されたBPMの5段階のレベル(1〜5)を格納\\ \hline
	  \texttt{currentVolStep} & \texttt{int}  & 算出された音量の5段階のレベル(1〜5)を格納 \\ \hline
\end{tabular}
\end{table}



\section{通信設計}
\subsection{通信の概要}
本設計では，Wifiを用いたUDPによる通信と，レーザーを用いた物理同期を行っている．通信の概要を図\ref{fig:gaiyou}に示す．
また，通信全体のシーケンス図を図\ref{zu:sikensu}に示す．
通信の手順としては，まず指揮デバイスでBPMまたは演奏状態の操作が行われた場合は，観覧車デバイスにUDPを用いて情報を送信する．
また，音量が操作された場合は，UDPを用いて楽器に情報を送信する．

次に，観覧車デバイスは指揮デバイスから送信されたBPMを基にモータを回転させる．
その回転周期を楽器デバイスのCdSセルで読み取り，BPMを算出する．

最後に指揮デバイスから各楽器に送る開始信号を送信し，それを受け取った楽器が演奏を開始する．
この時，BPMを基に開始信号を各楽器ごとにずらして送信することで輪唱を行う．
\begin{figure}[h]            
    \centering                
    \includegraphics[width=0.8\columnwidth]{./figs/sikensu.png} 
    \caption{通信全体のシーケンス図}
    \label{zu:sikensu}       
\end{figure}

\begin{figure}[H]
  \centering
  \includegraphics[width=0.8\linewidth]{./figs/gaiyou.pdf}
  \caption{本システムの通信経路の概要．}
  \label{fig:gaiyou}
\end{figure}
\subsection{ファイル構成}
各役割ごとに，ファイルを\ref{tab:commfile}に示すように定義する．
FerrisWheel.inoはメインの処理を行う．
MotorControl.cpp，MotorControl.hは観覧車の動作に関する処理を記述したファイルで，
SyncManager.cpp，SyncManager.hは指揮デバイスとの通信に関する処理を記述したファイルである．
\begin{table}[H]
\centering
\caption{観覧車に搭載するファイル構成．}
\label{tab:commfile}
	\begin{tabular}{l|p{22em}}
\hline
ファイル名 & 内容  \\
\hline
	FerrisWheel.ino & メインロジックを担当する．loop関数におけるUDP受信処理やモータ制御の呼び出しなどを行う．\\ \hline
	MotorControl.cpp & モーターの制御を行う．ルックアップテーブルに基づく回転速度の切り替えや，起動時の加速を担当する．\\ \hline
	MotorControl.h & MotorControl.cppのヘッダーファイルを定義する． \\ \hline
	SyncManager.cpp & 指揮デバイスとの通信を行う．UDPパケットの受信処理などを行う．\\ \hline
	SyncManager.h & SyncManager.cppのヘッダーファイルを定義する． \\ \hline
\end{tabular}
\end{table}


\subsection{通信に使用する変数・関数}
本設計に使用する定数は表\ref{tab:teisuu}の通りである．これらは実装において，プログラム内で const を使って定義する．
SSIDおよびPASSのメモリ量について，参考文献\cite{ref:Arduino_wifi}によると，Arduino WiFiライブラリではSSIDの最大長を32文字，Wi-Fiパスワード(WPAキー)の最大長を63文字として定義されている．また，Arduino言語を含むC言語やC++の仕様として，文字列の終点を表すヌル文字（\texttt{\textbackslash 0}）を格納するため1バイト分を加算する必要がある．そのため，本設計ではSSIDを33 byte，PASSを64 byteと算出した．
また，観覧車で使用する変数は表\ref{tab:hensuu1}の通りである．

次にArduinoのメモリの検証を行う．
本設計がターゲットハードウェアのメモリ制約を満たすか検証を行う．使用するマイコンボードはArduino UNO R4 WiFi である．データシートによれば，本機に搭
載されているマイコンは，プログラムの実行中に変数や計算結果を保持するSRAM を32 KB，プログラム本体を保存するFlash メモリを256 KB 備えている\cite{ref:Arduino_memori}．
表\ref{tab:teisuu}表および\ref{tab:hensuu1}に示した通り，本システムのメモリ使用量は 123 byte である．
これはSRAM 全体の約0.3\% 程度に留まっており，無線通信ライブラリや制御スタックによる動的なメモリ消費を考慮しても，十分な容量を確保できていると判断した．



本設計において，指揮デバイスと観覧車間の演奏制御および同期処理を実現するために実装する関数を表\ref{tab:function}に示す．

\begin{table}[H]
\centering
\caption{設計に利用する定数一覧．}
\label{tab:teisuu}
\begin{tabular}{lllll}
\hline
型 & 名前 & 説明 & 所属 & メモリ量 [byte] \\
\hline
float & VCC & Arduinoへの供給電圧$V_{cc}$ & Main & 4 \\
uint16\_t & LOCAL\_PORT & 通信に利用するポート番号 & SyncManager & 2 \\
char[] & SSID & Wi-Fiのネットワーク名 & SyncManager & Max:33 \\
char[] & PASS & そのWi-Fiのパスワード & SyncManager & Max:64 \\
uint8\_t & PIN\_MOTOR\_PWM & モーター制御用の出力ピン & MotorControl & 1 \\
uint8\_t & bpmTable[] & 5つの離散値BPM格納 & MotorControl & 5 \\
\hline
\end{tabular}
\end{table}


%さらに，bpmTable[]の中身は表\ref{tab:bpmtable}のように定義されている．
%\begin{table}[H]
%\centering
%\caption{bpmTable[] に格納するBPM値．}
%\label{tab:bpmtable}
%\begin{tabular}{c|ccccc}
%\hline
%インデックス番号 & 0 & 1 & 2 & 3 & 4 \\
%\hline
%BPM値 & 60 & 90 & 120 & 150 & 180 \\
%\hline
%\end{tabular}
%\end{table}


\begin{table}[H]
\centering
\caption{中継機で利用する変数一覧．}
\label{tab:hensuu1}
\begin{tabular}{lllll}
\hline
型 & 名前 & 説明 & 所属 & メモリ量 [byte] \\
\hline
uint8\_t & currentBPM & 現在演奏中のBPMデータ & MotorControl & 1 \\
float & RPM & BPMから計算された目標回転数 & MotorControl & 4 \\
float & omega & 目標角速度 $\omega$ & MotorControl & 4 \\
uint8\_t & pwmValue & analogWrite用の出力値 & MotorControl & 1 \\
float & V\_average & 平均電圧$V_{average}$ & MotorControl & 4 \\
\hline
\end{tabular}
\end{table}


\begin{table}[H]
\centering
\caption{実装する関数の一覧表．}
\label{tab:function}
\begin{tabular}{lllp{6cm}l}
\hline
戻り値 & 関数名 & 引数 & 概要 & 所属 \\ \hline
void & receiveUDP & なし & 指揮デバイスから送信されるBPMの段階番号を受信する．受信されたデータは currentBPM に格納する & SyncManager \\ \hline
void & startRamp & なし & 演奏開始時の加速を制御する．予約されたBPMに基づき1小節分の時間をかけて徐々に目標速度（RPM）へ到達させる． & MotorControl \\ \hline
void & changeBPM & uint8\_t & 演奏中のBPM変更を行う． & MotorControl \\ \hline
\end{tabular}
\end{table}

\subsubsection{クラス図}
本設計におけるクラス図を図\ref{fig:class}のように定義する．
メインとなる \texttt{Main} クラスが全体の進行を管理し，モータ制御を担う \texttt{MotorControl} クラスと通信・同期を担う \texttt{SyncManager} クラスを集約（所有）する構成とした．これにより，各機能の独立性を高めつつ，システム全体の司令を明確にしている．
\begin{figure}[H]
  \centering
  \includegraphics[width=0.8\linewidth]{./figs/classzu.png}
  \caption{本システムのクラス図．}
  \label{fig:class}
\end{figure}



\subsection{処理ロジック}
表\ref{tab:function}で定義した関数を用いて処理ロジックについて説明する．
無線通信は，利用するマイコンであるArduino UNO R4 WiFiに合わせるためWiFiS3.hライブラリを使用する．
指揮デバイス(指揮棒)を送信元とし，中継機へはBPMの段階番号を，楽器へはタイミングをずらしながら開始信号(ヘッダ)と音量をユニキャストで個別に送信する仕様とする．なお，終了時の演奏終了ヘッダは両方に送信する．
指揮デバイスに搭載された可変抵抗器（BPM，音量をユーザーが操作する用）および指揮棒(加速度センサを搭載した開始・終了検知用）の状態に基づき，表\ref{tab:range}のように電圧値から変換された1～5の段階番号や制御コマンドを受信する．
受信した段階番号を元に，中継機および楽器デバイスは表\ref{tab:bpm_volume}のような変換をする．

UDPで送信するパケットの構造は表\ref{tab:hed}，表\ref{tab:payload}のように，識別用のヘッダとデータ本体からなる合計2 バイトの固定長ペイロードとして定義する．
ここでUDPは到達保証がないプロトコルなので，システムの信頼性を高めるために冗長送信と送信待機を行う．
まず，冗長送信では，パケットロスによる制御の失敗を防ぐため，BPM変更，演奏開始・終了の各イベント発生時に，同一パケットを3回連続して送信する．
また，送信待機では指揮デバイスの状態が変化した時だけ，通信を行うことでネットワークの帯域の占有を最小限に抑える．



\begin{table}[H]
\centering
\caption{段階と入力値範囲の対応表．}
\label{tab:range}
\begin{tabular}{c|c}
\hline
段階 & 入力値の範囲 \\ \hline
1 & 0 -- 204 \\ \hline
2 & 205 -- 409 \\ \hline
3 & 410 -- 614 \\ \hline
4 & 615 -- 819 \\ \hline
5 & 820 -- 1023 \\ \hline
\end{tabular}
\end{table}

\begin{table}[H]
\centering
\caption{各段階に対応するBPMおよび音量設定値．}
\label{tab:bpm_volume}
\begin{tabular}{c|c|c}
\hline
段階 & BPM 設定値 & 音量設定値 \\ \hline
1 & 60 & 0 \\ \hline
2 & 90 & 64 \\ \hline
3 & 120 & 128 \\ \hline
4 & 150 & 192 \\ \hline
5 & 180 & 255 \\ \hline
\end{tabular}
\end{table}

\begin{table}[H]
\centering
\caption{通信の識別ヘッダ．(第1バイト:char)}
\label{tab:hed}
\begin{tabular}{c|c|l}
\hline
ヘッダ & 値（ASCIIコード） & 内容 \\ \hline
`'B'` & 0x42 & BPMデータ \\ \hline
`'V'` & 0x56 & 音量データ \\ \hline
`'S'` & 0x53 & 演奏開始（Start） \\ \hline
`'E'` & 0x45 & 演奏終了（End） \\ \hline
\end{tabular}
\end{table}

\begin{table}[H]
\centering
\caption{各ヘッダに対応したペイロードの内容．(第2バイト:uint8\_t)}
\label{tab:payload}
% 右側の列を l から p{22em} に変更し、自動改行させます
\begin{tabular}{c|p{22em}}
\hline
ヘッダ & ペイロード内容 \\ \hline
`'B'` & BPM の段階番号（1～5）を受信．中継機内部の \texttt{bpmTable[]} を参照し，設定値へ変換． \\ \hline
`'V'` & 音量 の段階番号（1～5）を受信．楽器ユニットへ送信される． \\ \hline
`'S'` & 演奏開始コマンド．各楽器ユニットが受信し，このパケットの到達をトリガーとして即座に演奏を開始． \\ \hline
`'E'` & 演奏終了コマンド．第2バイトにはダミーデータ（0）が格納されており，受信後は停止処理へ移行． \\ \hline
\end{tabular}
\end{table}



\subsection{回転速度同期機構の設計}
モータ制御の原理とBPMの変換について説明する．
BPMを，実際の回転数\texttt{RPM}に変換させることを目的とする．
本システムでは，ArduinoのPWM（パルス幅変調）機能を用いてモータの平均印加電圧を制御する．供給電圧を $V_{cc}$（定数 \texttt{VCC}），プログラム上の出力値を $pwmValue$ とすると，モータに加わる平均電圧 $V_{average}$ は(\ref{eq:average})式で決定される．
また，BPMをモータの目標回転数 \texttt{RPM} へ変換するロジックを定義する．本システムでは，4拍で90度回転する仕様に基づき，\texttt{RPM}は(\ref{eq:rpmteigi})式で表せ，1回転の時間は(\ref{eq:ssteigi})式で求められる．
制御の指標となる角速度\texttt{omega}($\omega$)は(\ref{eq:omegateigi})式で求められる．
\begin{equation}
\label{eq:average}
V_{average} = V_{cc} \times \frac{pwmValue}{256}
\end{equation}
\begin{equation}
\label{eq:rpmteigi}
  RPM = \frac{BPM}{16}
\end{equation}
\begin{equation}
\label{eq:ssteigi}
  60 \div RPM = 60 \div \frac{BPM}{16} = \frac{960}{BPM}
\end{equation}
\begin{equation}
\label{eq:omegateigi}
\omega = 360^\circ \div \frac{960}{BPM} = \frac{3}{8}BPM
\end{equation}

次にルックアップテーブルによる速度制御について説明する．
一般的なDCモータの回転速度は供給電圧に比例する性質があるが\cite{ref:control}，実際には摩擦や機構の負荷などによる非線形な挙動を示す事が多い．
そのため近似直線を用いる手法ではなく，より確実な速度追従を実現するため，実測値に基づくルックアップテーブル方式を採用する．

\texttt{MotorControl} クラスは，表\ref{tab:pwm_lut}に示す離散的な5段階のBPM設定値を保持し，受信した段階番号に応じた \texttt{pwmValue} を即座に選択・出力する構成とする．
\begin{table}[H]
\centering
\caption{BPMに対する目標RPMと実測出力PWM値}
\label{tab:pwm_lut}
\begin{tabular}{c|c|c|c}
\hline
段階 & 目標BPM & 目標RPM & 出力PWM値 \\ \hline
1 & 60 & 3.75 &（実測値A）\\
2 & 90 & 5.625 &（実測値B）\\
3 & 120 & 7.5 &（実測値C）\\
4 & 150 & 9.375 &（実測値D）\\
5 & 180 & 11.25 &（実測値E）\\ \hline
\end{tabular}
\end{table}
目標RPMは(\ref{eq:rpmteigi})式から算出している．ここからは実測値の測定方法について説明する．ただし，このロジックはシステムに本実装するプログラムではなく，測定方法であることに注意する．
まず，1象限分回転する時間$T_{target}$を(\ref{eq:target})式で計算する．
\begin{equation}
\label{eq:target}
    T_{target} = \frac{60}{BPM} \times 4
\end{equation}
そこから，任意のPWM 値でモーターを回し，実験用の外部計測器等を用いて任意のレーザーが1象限を通過する時間を\texttt{millis()}関数等を用いて測定する．この測定時間が$T_{target}$に限りなく近づくまでPWM 値を微調整する．

\subsection{レーザーによる物理同期信号の設計}
レーザーによるBPMの同期方法について説明する．
観覧車に45度間隔で搭載されたレーザーは常に点灯される．
観覧車の回転に伴って，レーザー光が楽器側のCdSセルを通過する際の光の立ち上がりを物理的な同期パルスとして利用する．
本設計では45度の移動を0.5 小節と定義しているため，このパルス検知は各小節の第1拍目と第3拍目の開始を意味している．

楽器ユニット側は，演奏開始前の空転時間において，このパルスの検出間隔から事前に実際のBPMを計算して演奏待機状態とする．実際の演奏開始は，指揮デバイスから個別に送信される演奏開始パケット（ヘッダ\texttt{'S'}'）を受信した瞬間をトリガーとする．

レーザー光によるBPMの同期処理は，通信クラスである \texttt{SyncManager} によって管理される．
具体的な処理としては，時間差によるBPM算出，BPMの補正，誤検知防止の3段階で行う．
まず，時間差によるBPM算出として観覧車のレーザー光がCdSセルを通過するたびに，
その時の時刻と前回の通過時刻との差分を測定する．
観覧車は，1小節が4拍とするとレーザーの間隔が0.5拍となり，回転速度から算出される
生のBPMは(\ref{eq:rawBPM})式で表される．
\begin{equation}
	\text{生のBPM} = \frac{30000}{\text{レーザー光検知時刻} - \text{前回の検知時刻}} \label{eq:rawBPM}
\end{equation}
次にBPM補正として，算出された生のBPMとあらかじめ決定していた演奏に使用するBPMのリストを比較し，
一番近いBPMを選択し，楽器へ送信する．
また，誤検知防止として，レーザー光がCdSセルの前を通過する際，細かな変動により一度の
通過で複数回検知される可能性を考慮して行う．
処理としては，レーザー光の検知直度に一定時間，検知をしないようにすることで複数回検知されることを防止し精度を高める．


%%%%%%%%%%%%%%%%%%%%%%%%
%%%%%%%%%%%%%%%%%%%%%%%%


\subsection{LEDマトリクスでのBPM表示}
まず，LEDマトリクスにおけるBPM表示について説明する．
今回のシステムでは表示用のマトリクスとして，Arduino Uno R4 WiFi付属のマトリクスを使用する．
このマトリクスを用いて，中継機から受信したBPM情報を元に文字列変換を行い，純正のパッケージを用いて表示を行う．
ここからは，本機能におけるファイル構成，関数設計について説明する．
\subsubsection{ファイル構成}
\begin{table}[h]
    \vspace{-1em}
    \caption{BPM表示機能のファイル構成}
    \centering
    \label{tab:file1}
    \begin{tabular}{|c|c|}
    \hline
    ファイル名 & 概要\\
    \hline
    Display.ino & 中継機からの通信管理\\
    \hline
    Display.cpp & 表示文字の独自フォント定義，LEDマトリクスへの描画処理\\
    \hline
    Display.h & 表示文字の独自フォント定義，LEDマトリクスへの描画処理\\
    \hline
    \end{tabular}
\end{table}
ファイル構成について説明する．また，表\ref{tab:file1}にBPM表示機能のファイル構成を示している．
はじめに，Display.inoファイルでは，表\ref{tab:file1}に示すように中継機からの通信管理を行う．
今回，中間機との通信形式としては，高速通信を目的とし，UDP形式で情報を受け取る．\par
次に，Display.cppおよび.hファイルでは，表示文字の独自フォント定義，マトリクスへの描画処理を行う．
ここで，今回利用するArduinoのマトリクスは縦12横8の範囲で構成されている．
この時，通常で定義されている文字列では文字の間隔が短くなることや，それに付随した可読性の低下が考えられる．
よって，今回は1文字を縦5横3で再定義することにより文字同士の間隔を開け，可読性の担保を目指す．
例として数字0を以下に示す．
\begin{verbatim}
{1,1,1,
1,0,1,
1,0,1,
1,0,1,
1,1,1}
\end{verbatim}
また，LEDマトリクスの描画処理は，受信したデータを文字列から整数に変換し，整数を各桁で分割，それぞれを定義した
配列に当てはめることで行う．描画自体はマトリクスパッケージの機能を用いて，配列を表現する．
\subsubsection{オブジェクト設計(クラス図)}
LEDMatrixプログラムのクラス図は以下の図\ref{matcla}の通りである．通信管理を行うDisplay.ino，LEDマトリクスへの描画処理を行うDisplay.cpp，および関数宣言やフォント定義を行うDisplay.hで構成される．
各モジュール間の関係について，Display.inoからDisplay.cppへの関係は，表示処理を利用する依存関係である．また，Display.cppからDisplay.hへの関係は，関数宣言およびフォント定義を参照する依存関係である．
\begin{figure}[h]
\centering
 \centering
\includegraphics[width=7cm]{./figs/class_mat.png}
\caption{LEDMatrixのクラス図}
\label{matcla}
\end{figure}

\subsubsection{関数設計}
表\ref{tab:func1}に今回用いる主な処理を行う関数を示している．以降では，各関数の具体的な動作を説明する．\par
はじめに，checkUDP関数について説明する．この関数は表\ref{tab:func1}に示すように，中継機からのパケット確認とBPM情報の値確認を行い，引数や返り値はない．
まず，中継機からのパケットをUdp.parsePacket()といったパケットを確認する関数を用いて監視する．もしパケットが届いていることが確認できれば，ペイロード（今回の場合はBPM情報）
を文字列として読み込む．その文字列に対し，規定のヘッダであるかを確認したのちに，BPM自体を数値として保持する．その後，値が更新されていれば，描画更新処理を行う．\par
次に，displayBPM関数について説明する．この関数は表\ref{tab:func1}に示すように，BPM数値のマトリクス描画を行い，引数としては，上記の保持した数値を用いる．
まず，現在のマトリクスを0でクリアし，表示をリセットする．その後，引数として保存したBPMの数値を1桁ずつに分割する．それらの値を定義してある独自フォント配列を照合し，
フレームバッファ内の指定位置にマトリクスの状態を描画する．そうして作成した配列データを，マトリクスパッケージに含まれるmatrix.renderBitmap()関数を用いてLEDとして発光させる．

\begin{table}[h]
    \vspace{-1em}
    \caption{主要な処理を行う関数一覧}
    \centering
    \label{tab:func1}
    \begin{tabular}{|c|c|c|c|}
    \hline
    関数名 & 概要 & 引数 &  返り値\\
    \hline
    checkUDP & 中継機からのパケットの確認，BPMの値確認 & なし & なし\\
    \hline
    displayBPM & BPM数値をマトリクスに描画 & int bpm & なし\\
    \hline
    \end{tabular}
\end{table}


\subsection{Tick管理およびシリアル通信}
%%%%%%%%%%%%%%%%%%%%%%%%%%%%%%%%%%%%%%%%%%%%%%%%%%%%%%%%%%%%%%%%%%%%%%%%%%%%%
演奏必要情報の生成処理について説明する．今回のシステムでは，中継機からBPMを受信し，それを元にして正確な時間であるTickを用いる．
本機能はこのTickの管理と，Processing側に対するシリアル通信を行うためのパケット通信を行う．
ここからは，本機能におけるファイル構成と関数設計について説明する．

\subsubsection{ファイル構成}
\begin{table}[bt]
    \vspace{-1em}
    \caption{楽器制御機能のファイル構成}
    \centering
    \label{tab:file2}
    \begin{tabular}{|c|c|}
    \hline
    ファイル名 & 概要\\
    \hline
    InstCtrl.ino & Tickの自律加算，processingへのシリアルパケット送信管理\\
    \hline
    Score.h & 楽譜データ構造体（NoteData）の定義\\
    \hline
    \end{tabular}
\end{table}

\begin{table}[bt]
    \vspace{-1em}
    \caption{楽譜データ構造体の定義}
    \centering
    \label{tab:deta}
    \begin{tabular}{|c|c|c|}
    \hline
    型 & 変数名 & 概要 \\
    \hline
    uint16\_t & timing & 発音タイミング（楽曲開始からの累積Tick数） \\
    \hline
    uint8\_t & pitch & 音高を示すMIDIノート番号（0〜127） \\
    \hline
    uint8\_t & startVelocity & 発音開始時の強さ（0〜127） \\
    \hline
    uint8\_t & endVelocity & 発音終了時の強さ（0〜127） \\
    \hline
    uint16\_t & duration & 発音の長さ（Tickまたはミリ秒） \\
    \hline
    \end{tabular}
\end{table}

ファイル構成について説明する．また，表\ref{tab:file2}に楽器制御機能のファイル構成を示している．\par
はじめに，InstCtrl.inoファイルではTickの自律加算，Processingへのシリアルパケット送信の管理を行う．
楽器用のArduinoでは，受信したBPMを用いてTickを算出し，その後は各楽器が内部のmicros()関数を用いてTickを加算し続ける．
ここで，中継機からBPMや音量の変化を受信した場合，発音より優先してそれらのシリアル送信を行う．
これは，ユーザー入力操作からのパラメータ変化を迅速に対応できるようにするためである．
それらが終われば，whileループを用いた楽譜走査のため，楽譜データの構造体から発音タイミングを示すtimingを読み込む．
そのパラメータを現在時刻と比較し，発音する音符を抽出する．抽出された音符はカンマ区切りの文字列１行（パケット）として保存され，
シリアル通信でProcessingへと送信される．\par
次に，Score.hでは，楽譜データ構造体（NoteData）の定義を行う．
また，表\ref{tab:deta}に今回定義する楽譜データ構造体の各変数名を示している．
楽譜データはArduinoのメモリ制限を考慮し，データサイズを明示的に指定できるstdint.hを用いる．
また，同様の理由により，楽譜データはconst型で定義し，フラッシュメモリ領域に保持する．

\subsubsection{関数設計}
\begin{table}[bt]
    \vspace{-1em}
    \caption{楽器制御機能の関数一覧}
    \centering
    \label{tab:func2}
    % 全体の幅を \textwidth にし、2列目と3列目を自動改行(X)に設定
    \begin{tabularx}{\textwidth}{|l|X|X|l|}
    \hline
    \textbf{関数名} & \textbf{概要} & \textbf{引数} & \textbf{返り値}\\
    \hline
    \texttt{updateTick} & 時間経過による自律Tick更新，センサによる同期修正 & なし & なし\\
    \hline
    \texttt{sendPacket} & 楽譜の走査，パケットの生成，送信 & なし & なし\\
    \hline
    \texttt{sendBpmValue} & BPM変更があった時の送信 & 読み取ったBPM & なし \\
    \hline
    \texttt{sendVolumeValue} & 音量の送信 & 受信した音量 & なし \\
    \hline
    \end{tabularx}
\end{table}

本機能の実装に用いる関数について説明する．また，表\ref{tab:func2}に今回用いる主な処理を行う関数を示している．以降では，各関数の具体的な動作を
説明する．\par
はじめに，updateTick関数について説明する．この関数は表\ref{tab:func2}に示すように，時間経過による自律Tick更新とセンサによる同期修正を行い，引数や返り値はない．
具体的には，中継機より受信したBPM情報から1Tickあたりの必要時間を算出する．また，センサ同期情報を受信し，それに合わせて
Tickの再算出を行い，可変のBPMに対応する．\par

次に，sendPacket関数について説明する．この関数は表\ref{tab:func2}に示すように，楽譜の走査とパケットの生成送信を行い，引数や返り値はない．
具体的には，上記の関数で計算したTickに基づいた速度で楽譜配列を走査し，発音タイミングとの比較を行ったのち，
パケットの生成とシリアル通信を行う．
シリアル通信の方法としては，write関数を用いて，あらかじめ定義したデータ構造をバイト単位の生データとして直接送信する．
具体的には．データの先頭を示す1バイトのヘッダーを送信した直後に，走査によって得られた「音の高さ」「発音時間」および「各種オプション」に対応するuint8\_t型のデータを順次書き出す．

sendBpmValueについて説明する．
この関数は，表\ref{tab:func2}に示すように，観覧車のLEDから算出したBPMが演奏中
のBPMと異なった時に，Processingへ送信する関数であり，引数は読み取った引数で戻り値はない．
具体的な送信手順として，まずsendPacket関数で送信するデータと違うことを示すためにBPM変更の識別用ヘッダーをwrite関数で送信する．
次に，uint8\_t型に格納したBPM値を1バイトで送信する．
このようにすることで，BPMの変更が起きた際のProcessingへの通信を行う．

sendVolumeValue関数について説明する．
この関数は，表\ref{tab:func2}に示すように，指揮デバイスから送信された音量を
Processingへ送信する関数であり，引数は音量で戻り値はない．
具体的な送信手順として，まず音量変更の識別のため，1バイトのヘッダーをwrite関数で送信する．
次に，uint8\_t型に格納した音量を1バイトで送信する．
これにより，音量をProcessingへ送信する．

%%%%%%%%%%%%%%%%%%%%%%%%%%%%%%%%%%%%%%%%%%%%%%%%%%%%%%%%%%%%%%%%%%%%%%%%%%%%%%%%%%%%%%%%%%%%%%%%

%%%%%%%%%%%%%%%%%%%%%%%%

%%%%%%%%%%%%%%%%%%%%%%%%

%%%%%%%%%%%%%%%%%%%%%%%%
%%%%%%%%%%%%%%%%%%%%%%%%
\section{処理フローの設計}
\subsection{全体のフローチャート}
今回製作するものの処理全体のフローチャートは以下の図\ref{flo}の通りである．
指揮デバイスの方で初期のBPMの設定を行った後，指揮棒を振った場合中継機デバイスや楽器に演奏情報を送信し，ボタンが押された際に通信開始情報を送信し，観覧車を作動させる．
\begin{figure}[H]
\centering
 \centering
\includegraphics[width=10cm]{./figs/flo.png}
\caption{全体のフローチャート図}
\label{flo}
\end{figure}



\subsection{LEDMatrixの処理フローチャート}
次に，LEDMatrixの処理フローについて，図\ref{flomat}を用いて説明する．
始めに初期化処理をしたあと，メインループに移行する．メインループは指揮棒から終了信号が来るまで行うループである．

メインループ内では，中継機から送信されるUDPパケットを常時監視する．パケットを受信した場合，文字列としてデータを取得し，ヘッダ確認を行う．ここで，ヘッダが規定値と一致しない場合は不正データとして破棄し，受信待機状態へ復帰する．
ヘッダが一致した場合はBPM値を取得し，前回値との比較によって値更新の有無を確認する．値が更新されていない場合は描画処理を行わず，受信待機状態へ戻る．一方，値更新が確認された場合はdisplayBPM関数を呼び出し，描画処理を実行する．
描画処理では，LEDマトリクスの表示状態を初期化した後，BPM値を各桁へ分割する．その後，各桁に対して独自フォント配列との照合および描画処理を繰り返し実行し，最終的にrenderBitmap()関数を用いてLEDマトリクスへ反映する．
その後，指揮棒が終了合図を送信し，それを受け取るとループから抜け出し，Matrix表示を終えると言う流れである．
\begin{figure}[H]
\centering
 \centering
\includegraphics[width=10cm]{./figs/flomat.png}
\caption{LEDMatrix内の処理フローチャート図}
\label{flomat}
\end{figure}

%%%%%%%%%%%%%%%%%%%%%%%%
%%%%%%%%%%%%%%%%%%%%%%%%
\section{テストの設計}
\subsection{音色の再現度テスト}
Processingで作成した音色(ピアノ，フルート，ストリングス，ドラム)の音と
，実際の楽器の録音音声を複数の被験者に聞いてもらい，似ているかの
5段階評価のアンケートに回答してもらう．
このアンケートの結果の中央値と最頻値を確認する．
\subsection{ボタン判定の制御}
指揮デバイスに接続されたボタンの制御プログラムのテストとして，
複数回ボタンを押下し，失敗する確率を測定する．
また，連打や長押しをした場合，意図しない動作をしないかを確認する．
成否の判定はシリアルモニタに結果を表示するプログラムを追加して行う．
\subsection{加速度センサの制御}
加速度センサの制御プログラムのテストとしては，まず適切な閾値accThresholdが正しく設定されているか確認する．
閾値が正しく設定されていないと，一度の振る動作で複数回検知されたり，振っても検知されないことが起きるためである．
具体的なテスト方法としてシリアルモニタに振ったと検知したタイミングで表示させるプログラムを追加し，適切に検知されるか確認する．
また，振ったと判断されるまでの遅延を計測し，体感として許容範囲であるかを評価する．

\subsection{BPM・音量読み取り機能}
BPM・音量読み取り機能のテストとして，まず移動平均のsampleSizeが適切かどうか
確認する．
sampleSizeが適切でないと，意図せずBPM・音量が変更されたり，移動平均算出にかかる時間が増加し，遅延が増加する可能性がある．
テスト方法としては，シリアルモニタに表示して意図せず変更されず，かつ遅延がどのくらいあるかを確認する．

\subsection{BPM・音量変更機能}
BPM・音量変更機能としては，正確に変更された情報が送信され受信できるかを確認する．
また，変更操作をしてから受信が完了するまでの時間を確認する．

\subsection{観覧車の制御}
観覧車の制御のテストとしては，設定されたBPMに合わせて4小節に正確に一回転しているかを確認する．また，最も低いBPM60の時モータが止まらずにスムーズな回転を維持できているかを確認する．
さらに，BPM60から180の急激なBPM変化で，スムーズに回転速度を変更できるかを確認する．
また，表示するLEDMatrixの内容が正しいことを確認する．


\subsection{BPM変更機能}
BPM変更機能のテストとしては，CdSセルがレーザー光の通過を正確に読み取れるかを確認する．
一度の通過で2回以上検知されたり，最高速度のBPM180でレーザー光を逃さずに検知できるかを確認する．
また，観覧車が動かしているBPMと読み取ったBPMが一致するかを確認する．


\subsection{同期精度テスト}
ここでは，物理的な動きと楽器の発音タイミングがどの程度一致しているかを検証することを目的とする．具体的には，観覧車のレーザー光が照度センサを通過する物理的なタイミングと，楽器側の発音タイミングが一致し，許容範囲内の遅延におさまっているかを確認する．

\subsection{負荷テスト}
ここでは，システムに高負荷をかけた際の挙動の安定性を検証する．
まず，通信負荷をテストするため，BPMや音量の変更パケットを短時間で連続して送信する．
次に物理負荷をテストするため，最大速度（180 BPM）で長時間連続運転を行う．
これらの操作により，システムが停止したり制御不能に陥ったりしないことを確認する．










\clearpage
\begin{thebibliography}{99}
\bibitem{cul} 株式会社マーケットリサーチセンター，``音楽出版の日本市場（～2031年）、市場規模（パフォーマンス、同期、デジタル収益）・分析レポートを発表,'' \url{https://www.atpress.ne.jp/news/4436457}，2026/4/22参照.
\bibitem{live} LP Information，``ライブコンサート市場占有率分析レポート,'' \url{https://www.atpress.ne.jp/news/3304316}，2026/4/22参照.
\bibitem{acpc} 一般社団法人コンサートプロモーターズ協会，``ライブ市場調査　基礎推移表'' \url{https://www.acpc.or.jp/marketing/transition/}，2026/5/18参照.
\bibitem{Natalia2018} Natalia Cooper, Dimitrios Mileounis, Patrick Williams, and Frank P. Brooks,`The effects of substitute multisensory feedback on task performance and the sense of presence in a virtual reality environment,''\textit{PLOS ONE},vol.13,no.2,p.e0191846,2018.
\bibitem{Martin}  Martin, R., and Nielsen, N. Enacting Musical Aesthetics: The Embodied Experience of Live Music. Music \& Science, 7,2024
\bibitem{yamaha} Yamaha Music Japan，``ENSPIRE PRO,'' \url{https://jp.yamaha.com/products/musical_instruments/pianos/disklavier/enspire_pro/features.html#product-tabs}，2026/5/19参照.
\bibitem{sekine2021} 関根 伸也，``DMX512信号システムなどの現状技術,'' 電気設備学会誌，vol.41，no.72026，p.382-385，2021
\bibitem{reactable} ReacTable Legacy，``Welcome|ReacTable Legacy,'' \url{https://reactable.com/}，2026/5/19参照.


\bibitem{bib:haas} H, Haas.: Über den Einfluss Eines Einfachechos auf die Hörsamkeit von Sprache, Acustica, Vol.1, No.2, pp.49-58 (1951). 
\bibitem{bib:twonine}KOBESOFT：稼働率，
	\url{https://kobesoft.co.jp/mikata/words/network-cloud/availability-rate/}，
		(2026/5/19参照)．
\bibitem{bib:tempo}
C. Drake and M.-C. Botte.:
Tempo sensitivity in auditory sequences: Evidence for a multiple-look model,
\textit{Perception \& Psychophysics},
		Vol. 54, No. 3, pp. 277-286 (1993).
	\bibitem{bib:ACSI} American Customer Satisfaction Index：U.S. Overall Customer Satisfaction，\url{https://theacsi.org/the-acsi-difference/us-overall-customer-satisfaction/}（参照2026-05-17）．

\bibitem{bib:tick} 有世犬, "ティック（tick） - DTM用語", \url{https://dtm.conceptmol.com/ティック（tick）/}, (参照 2026-05-22).

\bibitem{インハーモニシティ式}
中井 幹雄，川口 和也，
「弦のインハーモニシティに関する研究」，
日本機械学会，Dynamics \& Design Conference，2006，2006巻，セッションID109．

\bibitem{ref:Arduino_wifi} Arduino WiFi Library, wl\_definitions.h,, \url{https://github.com/arduino-libraries/WiFi/blob/master/src/utility/wl_definitions.h}, (参照 2026-05-08).

\bibitem{ref:Arduino_memori} Arduino, "Arduino UNO R4 WiFi Datasheet", \url{https://docs.arduino.cc/resources/datasheets/ABX00087-datasheet.pdf}, (参照 2026-04-28).

\bibitem{ref:control}
Control.com, ``DC Motor Speed Control,'' {\em Lessons In Electric Circuits}, \\
\url{https://control.com/textbook/variable-speed-motor-controls/dc-motor-speed-control/}，（参照 2026-05-04）．

\end{thebibliography}

\end{document}
